\documentclass[11pt]{article}
\usepackage[a4paper,margin=29mm]{geometry}
\usepackage{amsmath,amssymb,amsthm,mathtools,lmodern}
\usepackage[T1]{fontenc}
\DeclareMathSizes{10}{10}{8}{8}
\DeclareMathSizes{10.95}{10.95}{8}{8}
\usepackage[hidelinks]{hyperref}
\usepackage{microtype}
\setlength{\emergencystretch}{1.8em}
\newtheorem{proposition}{Proposition}[section]
\newtheorem{lemma}[proposition]{Lemma}
\newtheorem{theorem}[proposition]{Theorem}
\newtheorem{corollary}[proposition]{Corollary}
\newtheorem{claim}[proposition]{Claim}\renewcommand{\UrlFont}{\normalfont\ttfamily}
\newcommand{\Del}{\Delta}
\newcommand{\coef}[2]{[z^{#1}]#2}
\newcommand{\Iv}[2]{I_{[#1,#2]}}
\newcommand{\len}{\ell}
\renewcommand{\labelitemi}{\(\bullet\)}\renewcommand{\labelitemii}{\(-\)}\def\UrlBreaks{\do\_\do\.\do\/}
\newcommand{\wt}{\operatorname{wt}}
\newcommand{\pos}[1]{(#1)_+}


\title{Detailed coefficient lemmas}
\hypersetup{unicode=true,
  pdftitle={Detailed coefficient lemmas},
  pdfauthor={Zhang Yibo and Zhang Zhanhe}}
% The byline uses the title-page space otherwise reserved for an empty date.
\makeatletter
\renewcommand{\@maketitle}{%
  \begingroup\centering
  \vspace*{0.6em}%
  {\LARGE\@title\par}\vspace{0.85em}%
  \@author\par
  \endgroup
  \vspace{0.85em}}
\makeatother
\author{\begin{tabular}{@{}c@{\hspace{2em}}c@{}}
  \normalsize Zhang Yibo &
  \normalsize Zhang Zhanhe\\[-1pt]
  \footnotesize\href{https://orcid.org/0009-0002-2005-5781}{ORCID 0009-0002-2005-5781} &
  \footnotesize\href{https://orcid.org/0009-0007-2883-9640}{ORCID 0009-0007-2883-9640}
\end{tabular}\endgraf\vspace{3pt}%
{\footnotesize Both authors contributed equally as joint first authors. Zhang Zhanhe is the corresponding author.\endgraf}}
\date{}
\begin{document}
\maketitle
This appendix collects the mathematical source proofs used in the
strict-peak reduction of \emph{Rank-unimodality of one-descent composition
ideals}. The companion introduces their notation and gives the logical
order of use. Each excerpt identifies its source under \path{reduction/}.
The original files preserve the notation and the state of the investigation
at the time each lemma was proved. The single-arm inequality denoted by
(C) in these excerpts is Corollary 3.2 of the article.
\setcounter{tocdepth}{1}\tableofcontents\clearpage


\section{Ordinary products and the first mode}

\noindent{\small Source\enspace\path{formal_arm_extensions.tex}.}\par\medskip

\begin{proposition}[Ordinary products with no later return to one]
If $T=1^\alpha U$ and $W=1^\beta V$, with all parts of the
formal skeletons $U,V$ at least two, then $A_TA_W$ is unimodal.
The skeletons may have arbitrary internal order. Its reverse
also has unimodal positive first differences and is log-concave
on its rising side.
\end{proposition}
\begin{proof}
For every formal arm, the lowest height layer is $[0,a]$,
and every higher layer has endpoint sum at least $a+1$.
This follows by induction from $1+S_tA$: uniform convolution
gives all layers of $S_tA$ endpoint sum at least $a+t+1$,
and adding the initial constant only changes its lowest layer.
In the stated class, every nonrectangular polynomial has second
layer exactly $[2,a-1]$, since both coefficients at those two
ranks are at least two and the adjacent end coefficients are one.

The ordinary-product proof therefore applies verbatim within
this deletion-closed class. For completeness, write $a\ge b$.
Layer centers give growth through $\lceil(a+b)/2\rceil$.
When $a=b$ and both factors are nonrectangular, the product of
the two second layers pays for the only negative unit from the
two lowest layers at rank $a+1$, extending growth to that rank.
Deleting a last part $t$ of the heavier arm gives
$A_TA_W=A_W+S_t(A_{T^-}A_W)$. By induction the window term
is unimodal. If $a>b$, its tail joins the guaranteed increasing
segment after the correction $A_W$ has vanished; if $a=b$,
use the extra rank of growth just proved. Rectangular factors
are uniform-window base cases. Thus the product is unimodal.

The reverse-prefix proof in {\normalfont\path{joint_rising_differences.tex}}
uses only this product unimodality, the same strengthened midpoint,
and the formal append recurrence. Those inputs all hold within
the present deletion-closed class, proving the two additional
reverse properties.
\end{proof}

\noindent{\small Source\enspace\path{joint_rising_differences.tex}.}\par\medskip

\begin{proposition}
Let $X,Y$ be increasing arms and
$p(z)=z^{a+b}F_X(1/z)F_Y(1/z)$, with $a=\operatorname{wt}X$
and $b=\operatorname{wt}Y$. Then $p$ is unimodal and log-concave
on its rising side. In addition, the positive part of its
first-difference sequence is unimodal.
\end{proposition}
\begin{proof}
Ordinary two-arm product unimodality is an established input;
the proof by interval layers and final-part deletion is reproduced
in the supplied foundation. Its strengthened midpoint assertion
says that, when $a=b$ and both factors are nonrectangular,
$(F_XF_Y)_{a+1}\ge(F_XF_Y)_a$.

Induct on the total number of arm parts for the two additional
properties. If either arm polynomial is an interval, apply the
uniform-window results to the reverse of the other arm. Otherwise
orient $a\ge b$ and write $X=Zt$. Reversal gives
\[
 p=[t]q_Zq_Y+z^a q_Y.
\]
Let $m_0$ be the first mode of $p$. First-half increase of the
ordinary product and its unimodality imply that its last mode
is at least $\lceil(a+b)/2\rceil$. Thus
$m_0\le\lfloor(a+b)/2\rfloor$. If $a>b$, this gives $m_0<a$.
If $a=b$, the strengthened midpoint assertion implies that the
last forward mode is at least $a+1$, and hence $m_0\le a-1$.
This reasoning explicitly allows a maximal plateau.

In both weight cases $m_0<a$. Up to and including
$m_0$, $p$ agrees with the uniform window $[t]q_Zq_Y$; the
added term begins at $a$. The positive first differences of
$p$ are therefore a prefix of those of this window, followed
by zeros. The window has unimodal positive first differences
by induction and the uniform-window lemma. This proves that
the positive first differences of $p$ are unimodal.

For log-concavity before the first mode, the same agreement
and the rising-side window lemma apply. A window cannot have
passed its first mode and then strictly increase to the first
mode of $p$; consequently its relevant prefix lies on its
rising side. Log-concavity at a mode is automatic, since both
neighbors are at most the mode value. Interior points of a
maximal plateau also satisfy equality. This proves rising-side
log-concavity of $p$, including its last mode.
\end{proof}


\section{The two weighted base cases}

\noindent{\small Source\enspace\path{binary_ternary_A2_original_family_proof.tex}.}\par\medskip

First suppose one arm is rectangular, $q_X=[\alpha+1]$, and put
$b=\deg q_Y$. If $\alpha\ge b$, the first ordinary product mode is
$b$. Its positive differences are $(q_Y)_0,\ldots,(q_Y)_b$, then
zeros. Their weighted values increase, since actual ideal growth
gives $2(q_Y)_n\ge(q_Y)_{n-1}$. If $\alpha<b$, the first product
mode $\sigma$ is below $b$: at rank $b$ the window difference is
$1-(q_Y)_{b-\alpha-1}\le0$, and ordinary window unimodality applies.
The single-arm corrected source
$e_Y=(2-z)q_Y-z^b+z^{b+1}$ is nonnegative unimodal.
The product source equals $[\alpha+1]e_Y$ at every rank below $b$.
Its prefix through $\sigma$ is unimodal, and (3), followed by the
zero positive-difference tail, proves A2. This includes the empty
arm and two rectangles. We do NOT assert full raw-source shape
for rectangular pairs.

\noindent{\small Source\enspace\path{deficit_cumulative_ratios.tex}.}\par\medskip

For an increasing arm $T$ of weight $c$, write
\[
 E_T=1-z+(2z-1)F_T,\quad R_T=z^{c+1}E_T(1/z),\quad
 C_T(n)=\sum_{j=0}^n[z^j]R_T.
\]
All cumulative sequences are zero at negative indices and constant
after their finite input ends. The empty arm has $R_\varnothing=1$.

\begin{theorem}\label{thm:deficit-cumulative-lc}
For every increasing arm, $C_T$ is log-concave on all integer indices.
\end{theorem}
\begin{proof}
Induct using $R_{Tt}=[t]R_T+z^c$. The empty case is constant one
on the nonnegative integers. For a nonempty prefix the positive
support of $R_T$ is $[0,d]$, where $d=c+1-\kappa\le c-1$;
the increasing condition gives $\kappa\le t+1$. Its mass $m$ equals
$F_T(1)\ge c+1$, since every rank from zero to $c$ is nonempty.

Set $w=[t]R_T$ and $v_n=\sum_{j=0}^n w_j=[t]C_T(n)$.
By induction and log-concave convolution, $v$ is log-concave.
The new cumulative sequence is $h_n=v_n+\mathbf1_{n\ge c}$.
We check the joining edge and subsequent centers.

If $c\ge t$, then
\[
 w_{c-1}-w_c=R_T(c-t)\ge1.
\]
Indeed $0\le c-t\le d$. The known unimodality of $R_T$
and of a uniform window makes $w$ nonincreasing from $c-1$
onward. Thus $v$ is concave there. Writing
$V=v_{c-2}$, $x=w_{c-1}$ and $y=w_c$, the determinant at
the only left joining center is
\[
 h_{c-1}^2-h_{c-2}h_c=V(x-y-1)+x^2\ge0.
\]

If $1\le c<t$, then $w_n=m$ on $[c-1,t-1]$ and is
nonincreasing afterward. Again $v$ is concave from $c-1$.
At the joining center its determinant is
\[
 m^2-v_{c-2}\ge m^2-(c-1)m\ge0.
\]
In both cases the determinant at $c$ is the old determinant
plus $2v_c+1-v_{c-1}\ge0$. At every center $n\ge c+1$ it is
the old determinant plus $2v_n-v_{n-1}-v_{n+1}\ge0$.
Earlier centers are unchanged. Finally, for the empty prefix
$c=0$, the new cumulative sequence is $2,3,\ldots,t+1$ followed
by a constant tail, which is log-concave. This closes induction.
\end{proof}

\begin{lemma}[A cumulative premise implies actual prefix-to-tail ratios]
Let $p$ be a finite nonnegative sequence, $p_0>0$, and suppose
its cumulative sequence $C_n$ is log-concave. For every $t\ge2$,
the actual coefficients $w=[t]p$ satisfy
\[
 w_{k+1}w_j\ge w_kw_{j+1}\qquad(0\le k\le t-2,\ j\ge k).
\]
\end{lemma}
\begin{proof}
Fix $k$ and put $r=C_{k+1}/C_k$. Since $w_k=C_k$ and
$w_{k+1}=C_{k+1}$, it suffices to show $w_{j+1}\le rw_j$.
For $k\le j\le k+t$, cumulative log-concavity gives
$C_{j+1}\le rC_j$ and $C_{j+1-t}\ge rC_{j-t}$;
zero extension covers negative lower indices. Subtraction proves
the claim. For $i\ge k+1$, the same premise gives
\[
 p_{i+1}C_{i-1}\le p_iC_i,\qquad
 p_{i+1}\le (C_i/C_{i-1})p_i\le rp_i.
\]
Summing these inequalities over the window proves the remaining
case $j\ge k+t$. All denominators used are positive.
\end{proof}

\begin{corollary}\label{cor:deficit-window-ratios}
For every increasing arm $T$ and every $t\ge2$, $[t]R_T$
has the prefix-to-tail coefficient comparison through $t-2$.
The window need not exceed the largest part of $T$.
\end{corollary}

\begin{proposition}[The heavier one-part terminal case at weight two]
Let $A=F_{1^\alpha t}$ and let $B$ be any nonrectangular
increasing-arm polynomial. Suppose $t\ge2$, $\alpha\ge1$,
and $a=\alpha+t\ge b=\deg B$. For $q=z^{a+b}(AB)(1/z)$,
the sequence $2^j\max(q_j-q_{j-1},0)$ is unimodal.
\end{proposition}
\begin{proof}
The identity $A=1+z[t][\alpha+1]$ permits exchange of the
two interval factors. We may therefore assume $\alpha\ge t-1$.
The exchanged capacity is smaller if an exchange was needed.
Put $q_B=z^bB(1/z)$ and $V=[t]q_B$. Then
\[
 q=[\alpha+1]V+z^a q_B.
\]
Its first mode $s$ is strictly less than $a$, by ordinary-product
first-half growth and the extra midpoint rank for equal nonrectangular
arms. Through $s$, its positive differences are
$d_j=V_j-V_{j-\alpha-1}$. They are $V_j$ for $j\le\alpha$.
The sequence $2^jV_j$ is nondecreasing: in forward order the
coefficients of $[t]B$ grow by a factor at most two. For ranks
below $t$ use $f_n\le\sum_{i<n}f_i$; afterward use
$f_n\le2f_{n-1}$ termwise.

For $\alpha+1\le j\le s$, set $k=j-\alpha-1\le t-2$.
Also $s\le\deg V=b+t-1$, so $k<b$ by normalization.
Writing $W=[t]R_B$, the exact identity is
\[
 (2-z)V=W+z^b-z^{b+t},\qquad
 2d_j-d_{j-1}=W_j-W_k+\mathbf1_{j=b}.
\]
The negative terminal monomial is outside $j\le s$ and neither
monomial affects the head index $k<b$. By the preceding corollary,
$W_{k+\alpha+1}/W_k$ is nonincreasing wherever its denominator
is positive. Thus $W_j-W_k$ has no negative-to-positive change.
If such a difference was already negative before $j=b$, it is
at most $-1$ at $b$, by integrality and the ordered ratios; the
extra unit still cannot make it positive. Zero endpoints are
handled directly. The actual weighted differences have the same
single-crossing property. Joining to their nondecreasing initial
segment and then to terminal zeros proves the proposition.
\end{proof}

\section{Tilt and the common rising prefix}

\noindent{\small Source\enspace\path{clipped_forward_ratio_frontier.tex}.}\par\medskip

Let $A=F_X$, $B=F_Y$, with weights $a,b$, and put
\[
 D_n=\max(A_n-A_{n-1},0)\quad(n\ge2),\qquad H=E_B A.
\]
Here $D_0=D_1=0$. The following is an \emph{unproved}
generation-dependent comparison:
\[
 \frac{\min(H_n,D_n)}{D_n}\le
 \frac{\min(H_{n+1},D_{n+1})}{D_{n+1}}
 \quad\left(2\le n<\min(a,b),\ D_nD_{n+1}>0\right).
                                                        \tag{R}
\]
All ratios can be cross multiplied. Full, unclipped ratio
monotonicity is false; subtracting a single unit does not
repair it either. The clipping and the stated range must
not be silently removed.

\begin{proposition}[Conditional coefficient-tilt consequence]
If (R) holds for a given pair of increasing nonrectangular
arms, then the coefficients of $A(z/2)B(z/2)$ are unimodal.
Equivalently $2^j(q_Xq_Y)_j$ is unimodal.
This is coefficient tilt, not the weighted-positive-difference
property needed to finish the original conjecture.
\end{proposition}
\begin{proof}
The exact identity
\[
 (2z-1)AB=E_B A-(1-z)A                              \tag{1}
\]
shows that, while $A$ is rising, its rank-$n$ coefficient
is $H_n-D_n$. Once $H_n\ge D_n>0$, (R) forces
$H_{n+1}\ge D_{n+1}$ whenever the latter denominator is
positive. If it vanishes, initial interval support of the
ordinary positive differences shows that $A$ has finished
rising; the coefficient in (1) is then nonnegative forever.

Only ranks through $\min(a,b)$ require this argument.
Orient $a\le b$. For $n>a$, every nonzero summand in
$[z^n]AB=\sum_{i=0}^{a}A_iB_{n-i}$ uses $n-i\ge1$.
The elementary actual-ideal growth bound
$B_r\le2B_{r-1}$ therefore gives
$[z^n]AB\le2[z^{n-1}]AB$ directly. The initial ranks
are $1,2,5$ in $AB$, so its tilt starts with a plateau
and a rise. Identity (1), (R) and the terminal growth bound
exclude every later decrease followed by a strict rise.
\end{proof}


\noindent{\small Source\enspace\path{clipped_forward_ratio_frontier.tex}.}\par\medskip

For a single increasing nonrectangular arm $A$, write
$G_j=\sum_{i\le j}A_i$, zero when $j<0$. The following
single-arm assertion is still \emph{unproved}:
\[
 G_{n+1-r}D_n\ge G_{n-r}D_{n+1}
 \quad\left(r\ge3,\ D_nD_{n+1}>0,\ n\ge2\right).
                                                        \tag{C}
\]
Unlike (R), this assertion has no clipping. Broad and targeted
finite tests do not constitute its proof.

\begin{proposition}[Conditional single-arm reduction]
If (C) holds for every increasing arm, then (R) holds for
every pair of such nonrectangular arms. Consequently their
ordinary product has unimodal factor-two coefficient tilt.
\end{proposition}
\begin{proof}
At a pair of ranks where $B$ is still strictly rising,
the established single-arm deficit mode ordering gives
$\Delta E_B(j)\ge0$ throughout the prefix in question.
Since $E_B(j)=0$ for $j\le2$, summation of its differences
gives the exact coefficient identity
\[
 H_n=(E_B A)_n
       =\sum_{j=3}^{n}\Delta E_B(j)\,G_{n-j}.
                                                        \tag{2}
\]
Apply (C) to every $j$ in this sum. The new term at $j=n+1$
has zero old coefficient and a nonnegative new coefficient.
Multiplying by the nonnegative $\Delta E_B(j)$ and summing
therefore proves the \emph{unclipped} ratio comparison
$H_{n+1}D_n\ge H_nD_{n+1}$ in this range. Clipping preserves
this comparison at threshold one.

If $B$ has already ceased strictly rising at rank $n+1$,
its positive-difference initial support implies it never
strictly rises again. Use the symmetric version of (1):
\[
 (2z-1)AB=E_A B-(1-z)B.
\]
Its coefficients at this and all later ranks are nonnegative,
because $E_A B\ge0$ and $B_j-B_{j-1}\le0$ there.
Thus $H_{n+1}\ge D_{n+1}$ whenever $D_{n+1}>0$, and
the new clipped ratio is one. The same argument handles
leaving the rising support of $A$ as in the earlier proof.
These cases establish (R).
\end{proof}

\noindent{\small Source\enspace\path{C_to_coupled_S_rising_window.tex}.}\par\medskip

\begin{theorem}[Nonrecovering reverse deficit from one-arm C]
Let $A=F_Z$, $B=F_Y$ be actual increasing nonrectangular arms, and
assume $A$ has property C. Set $D_i=(A_i-A_{i-1})^+$,
$D_0=D_1=0$, $E_B=1-z+(2z-1)B$, and $H=E_BA$.
At a forward comparison $j\ge2$, assume
$D_jD_{j+1}>0$ and $E_B$ is nondecreasing throughout ranks
zero through $j+1$. If
\[
                    2(AB)_j-(AB)_{j+1}\ge0,
\]
then
\[
            2(AB)_{j-1}-(AB)_j
                       \le2(AB)_j-(AB)_{j+1}.           \tag{1}
\]
In particular the deficit monotonicity premise holds whenever $B$
is still strictly rising at $j+1$, by the established single-arm
ordinary/deficit mode ordering. If both ordinary factors are still
strictly rising, the theorem applies whenever either one has C.
\end{theorem}
\begin{proof}
Write $w_r=E_B(r)-E_B(r-1)$. Nonrectangularity gives $w_r=0$
for $r<3$; the prefix hypothesis gives $w_r\ge0$ through $j+1$.
The exact source identities are
\[
 H_j=\sum_{r=3}^j w_rG_A(j-r),\qquad
 H_{j+1}=\sum_{r=3}^{j+1}w_rG_A(j+1-r).
\]
Multiply the shift-$r$ C inequalities by $w_r$ and sum, retaining
the nonnegative newly born $r=j+1$ term. This proves
$H_{j+1}D_j\ge H_jD_{j+1}$. Also
$H_{j+1}-H_j=\sum_rw_rA_{j+1-r}\ge0$.

The exact identity
\[
                    (2z-1)AB=E_BA-(1-z)A
\]
identifies the two deficits in (1) with $H_j-D_j$ and
$H_{j+1}-D_{j+1}$. The latter is nonnegative by hypothesis.
These facts imply its monotonicity by the positive decomposition
\[
\begin{split}
 H_{j+1}\bigl((H_{j+1}-D_{j+1})-(H_j-D_j)\bigr)
  ={}&(H_{j+1}D_j-H_jD_{j+1})\\
    &+(H_{j+1}-H_j)(H_{j+1}-D_{j+1})\ge0.
\end{split}
\]
Here $H_{j+1}\ge D_{j+1}>0$, so division proves (1).
This step uses BOTH the unclipped ratio comparison and monotonicity
of the actual source convolution; clipping alone would not suffice.
\end{proof}

\begin{corollary}[Correct coupled S on the two-rising-factor branch]
Put $p=q_Zq_Y$, $a=\operatorname{wt}Z$, $b=\operatorname{wt}Y$,
and $R_i=2p_i-p_{i-1}$. At a usual S crossing index $n$ set
$j=a+b-n$. If $j\ge2$, both $F_Z,F_Y$ are still strictly rising
at $j+1$, at least one of them has C, and $R_n\ge0$, then
\[
                         R_{n+1}\le R_n.                \tag{2}
\]
Assume, as in the A2 induction, that the OLD weighted positive
differences $2^iu_i$ are unimodal. For $k<\ell$, with $\ell$ their
first mode, one has $u_k\le2u_{k+1}$. Formula (2) is exactly
$v_n\le2v_{n+1}$, where $v_n=p_{n-1}-p_n$ in the crossing
range. Adding gives
\[
                    u_k+v_n\le2(u_{k+1}+v_{n+1}),
\]
which is S. Strict weighted fall is not needed inside this positive
branch, although it cannot be removed from general S.
\end{corollary}

\section{The common nonincreasing phase}

\noindent{\small Source\enspace\path{agent_bothdeclining_first_birth.tex}.}\par\medskip

Let $Z,Y$ be actual increasing normalized arms, of weights $a,b$,
with forward polynomials $A,B$, and write $P=AB$,
$q_Z=z^aA(1/z)$, $q_Y=z^bB(1/z)$.
The POSITIVE corrected reverse source is
\[
 \widetilde e_Y=z^{b+1}E_Y(1/z)
    =(2-z)q_Y-z^b+z^{b+1},
 \qquad E_Y=1-z+(2z-1)B.
\]
It must not be confused with the raw polynomial $(2-z)q_Y$.
For every actual increasing arm, $E_Y$ and $\widetilde e_Y$ are
nonnegative and unimodal by the legal-maximum append argument in
{\normalfont\path{deficit_unimodality.tex}}. Every height interval of
$\widetilde e_Y$ has endpoint sum at most $b-1$, by
{\normalfont\path{deficit_forward_reflection.tex}}. Every height interval
of $q_Z$ has endpoint sum at most its degree $a$.
These are established SINGLE-arm inputs, independent of general A2.

Put $\kappa=\kappa_Y$, $r=a-j$, $n=a+b-j$, and
\[
 C_Z=\Delta q_Z(r+1)-\Delta q_Z(r),\qquad
 T=P_{j+1}-3P_j+2P_{j-1}.
\]
All coefficient sequences are zero-extended.

\begin{proposition}[Unconditional first-height payment]
If $j\le b$ and $2j\le a+b+2$, then
\[
 \boxed{T\le C_Z-q_Z(r+\kappa-1).}                 \tag{1}
\]
In particular (1) holds whenever $j\le\min(a,b)$.
No source-prefix, OLD A2, weighted-tail, head, or mode assumption
is used in this assertion.
\end{proposition}
\begin{proof}
The corrected source identity is
\[
 (2-z)q_Zq_Y=\widetilde e_Yq_Z+z^b(1-z)q_Z.
\]
Taking its difference at $n+1$ gives exactly
\[
 T=C_Z-\{(\widetilde e_Yq_Z)_n
                  -(\widetilde e_Yq_Z)_{n+1}\}.    \tag{2}
\]
Every product of one height interval of $\widetilde e_Y$ and
one height interval of $q_Z$ is symmetric unimodal and has
endpoint sum at most $a+b-1$. The assumption
$2j\le a+b+2$ says $2n+1\ge a+b-1$, so every such
component has a nonnegative drop from $n$ to $n+1$.

The lowest height interval of $\widetilde e_Y$ is exactly
$[0,b+1-\kappa]$: the forward source has positive interval
support $[\kappa,b+1]$, and its final coefficient is two.
The drop contributed by this ONE layer convolved with $q_Z$ is
\[
 q_Z(n-(b+1-\kappa))-q_Z(n+1)
   =q_Z(r+\kappa-1),
\]
because $j\le b$ implies $n\ge a$ and hence $q_Z(n+1)=0$.
All remaining layer-product drops are nonnegative. Substitution
in (2) proves (1).
\end{proof}

\paragraph{Consequences of the already proved first-birth payments.}
Suppose $Z=1^\alpha U$, $h=|U|\ge2$, $D=a-\alpha-h$,
and $A_{j+1}\le A_j$. Throughout the FULL region
$j\le b$, $2j\le a+b+2$,
all first-birth arguments in {\normalfont\path{source_birth_gap_nonrecovery.tex}},
{\normalfont\path{source_birth_gap_refinements.tex}}, and
{\normalfont\path{arbitrary_source_gap_finite_owner.tex}} remain valid
with (1) replacing the nondecreasing-source convolution bound.
Their scalar/primitive/reflection statements use no other source-prefix
assumption. Consequently:
\begin{itemize}
\item $\alpha\ge\kappa-5$ forces $T\le0$;
\item if $h\ge4$, $\alpha\ge\kappa-6$ forces $T\le0$;
\item if $T>0$, then, for $d=\kappa-\alpha$,
      $d\ge6$ and $D\le2d(d-3)$, hence
      $a\le\alpha+4d(d-3)$ and $j\le a+1$.
\end{itemize}
The earlier gap-five and gap-six proofs retain their exhaustive finite
small-reverse-rank complements; this note supplies a new source payment,
not a new independent formalization of those complements.
The negative reverse-rank cases also transfer: $r\le-2$ has
$C_Z=0$, and $r=-1$ has $C_Z=1$, paid by the same support
coefficient used in the earlier birth arguments. Thus the displayed
full region is not being silently restricted to $r\ge0$.
If BOTH $A$ and $B$ are nonincreasing and $j\le\min(a,b)$, these
necessary bounds apply in BOTH orientations. This does not itself
exclude simultaneous first-birth failure; the capacity argument below
does exclude it in this common-short-rank region.

\begin{proposition}[The second lowest layer pays the second birth]
Assume the first-height proposition's full scope, $\kappa\ge4$,
and $A_{j+1}\le A_j$. Put $v=r+\kappa-1$. If
\[
                         C_Z\le q_Z(v)+q_Z(v+2),     \tag{3}
\]
then $T\le0$.
\end{proposition}
\begin{proof}
The early normalized source bound gives $E_Y(\kappa+2)\ge2$.
Since $\widetilde e_Y(0)=2$, the SECOND height interval also
starts at zero. Its right endpoint is $b+1-\kappa-\delta$
for some $\delta\in\{0,1,2\}$. At the compared edge its drop
after convolution with $q_Z$ is exactly $q_Z(v+\delta)$.
The first two height layers are distinct, even when their
endpoints agree, so
\[
                T\le C_Z-q_Z(v)-q_Z(v+\delta).
\]
If $C_Z\le q_Z(v)$, the first-height proposition already pays.
Otherwise $C_Z>q_Z(v)\ge0$, so $r\ge-1$; the premise on $A$
gives $\Delta q_Z(r)\ge0$, including its zero extension at $r=-1$.
Hence
\[
 q_Z(r+1)\ge\Delta q_Z(r+1)\ge C_Z>q_Z(v),
 \qquad r+1<v.
\]
Unimodality of $q_Z$ therefore puts $v$ after its first mode
and makes $q_Z(v+\delta)\ge q_Z(v+2)$.
Now (3) pays the two-layer bound.
\end{proof}

Consequently the full absolute first-weight-eight theorem for
$h\ge4$ from {\normalfont\path{source_birth_gap_refinements.tex}} also
transfers to $j\le b$, $2j\le a+b+2$ without a source-prefix
assumption. Its only non-first-birth finite complement uses exactly
the two-term criterion (3); the same canonical initial exchange,
exhaustive small-rank complement, and earlier first-birth proof apply.
Thus an owner with at least four non-one parts has $T\le0$
throughout this region whenever $\kappa_Y\le8$.

\begin{proposition}[A later source increment with its subtraction removed]
Assume $2j\le a+b+2$, let $L\ge2$ and $L\le\kappa-3$, and suppose
\[
                       b-j\ge\kappa+L-2.            \tag{4}
\]
Then
\[
 \boxed{T\le C_Z-2^{L-2}q_Z(r+\kappa-1+L).}         \tag{5}
\]
\end{proposition}
\begin{proof}
The proved early source-curvature profile gives
$\Delta E_Y(\kappa+L)\ge2^{L-2}$.
Under reversal, this counts at least $2^{L-2}$ distinct
height intervals of $\widetilde e_Y$ whose RIGHT endpoint is
$R=b+1-\kappa-L$. For each such interval $[\ell,R]$,
its endpoint-sum bound gives
\[
 \ell+R\le b-1,\qquad \ell\le\kappa+L-2\le b-j.
\]
Its drop after convolution with $q_Z$ is
\[
 q_Z(n-R)-q_Z(n+1-\ell)
    =q_Z(r+\kappa-1+L),
\]
since $n+1-\ell\ge a+1$ by (4). Every unselected
layer-product has a nonnegative drop by the same upper-center
comparison as before. Equation (2) proves (4).
\end{proof}

Thus the linear-degree argument of
{\normalfont\path{source_curvature_linear_owner.tex}} also transfers
to this region whenever its $L=d-(D-2r-2)-3\ge2$ satisfies (4).
For $L=0,1$, that proof uses only the first birth, so (4) is unnecessary.
Without (4), dropping the second coefficient of a higher height interval
would be unjustified. In particular no global source-prefix assumption
has silently been replaced by an unconditional positive convolution.

\begin{lemma}[The owner's last capacity pays a primitive curvature]
Suppose $Z=Z^-t$, $t\ge2$, and $A_{j+1}\le A_j$.
Write $r=a-j$ as above. For every integer $1\le\delta\le t$,
\[
                    C_Z\le q_Z(r+\delta)
                 \quad\hbox{whenever }C_Z>0.        \tag{6}
\]
\end{lemma}
\begin{proof}
For $r\le-2$ the premise $C_Z>0$ is impossible. For $r=-1$,
$C_Z=1$ and $q_Z(\delta-1)\ge1$, since
$0\le\delta-1\le t-1<a$. Assume $r\ge0$.
The last two ordinary coefficients of $q_Z$ are both one.
Thus $\Delta q_Z(r)\ge0$ and $C_Z>0$ force $r\le a-2$:
at $r=a-1$, the first condition gives $q_Z(a-2)\le1$, whereas
the second would require $q_Z(a-2)>1$; the later cases have
nonpositive curvature or a negative first difference.

Put $Q=q_{Z^-}$ and $W=[t]Q$. The terminal unit of
$q_Z=W+z^a$ is absent through $r+1$, so
\[
 \Delta q_Z(r)=Q_r-Q_{r-t}\ge0,\qquad
 C_Z=\Delta Q(r+1)-\Delta Q(r+1-t)>0.                \tag{7}
\]
If $\Delta Q(r+1)\le0$, (7) gives
$\Delta Q(r+1-t)<0$. Unimodality of $Q$ then makes it
nonincreasing from this strict negative edge onwards. Hence
$Q_r\le Q_{r+1-t}<Q_{r-t}$, contradicting (7).
Therefore $\Delta Q(r+1)>0$. No earlier difference of a
unimodal sequence can be negative before a positive difference,
so $\Delta Q(r+1-t)\ge0$. Consequently
\[
 C_Z\le\Delta Q(r+1)\le Q_{r+1}.
\]
The interval sum $W_{r+\delta}$ includes the summand $Q_{r+1}$,
because its contributing old indices range from
$r+\delta-t+1$ to $r+\delta$ and $1\le\delta\le t$.
All summands are nonnegative, proving (6).
\end{proof}

\begin{theorem}[Complete common-short-rank both-nonincreasing theorem]
Suppose BOTH arms are nonrectangular, BOTH ordinary factors
are nonincreasing at $j+1$, and
$j\le\min(a,b)$. Then $T\le0$, with no source-prefix,
nonnegative-tail, OLD A2, head, or product-mode assumption.
\end{theorem}
\begin{proof}
Orient the arms so the last capacity $t$ of $Z$ is at least
the last capacity of $Y$. The first missing weight satisfies
$\kappa_Y\le\max(Y)+1$, so $1\le\kappa_Y-1\le t$.
If $C_Z\le0$, the first-height proposition pays immediately.
Otherwise apply (6) with $\delta=\kappa_Y-1$, and then (1).
\end{proof}

More generally, within the first-height proposition's full scope,
a nonincreasing owner $Z$ with last capacity at least
$\kappa_Y-1$ always has $T\le0$. In particular a positive
recovery in that full scope requires
$\kappa_Y\ge\max(Z)+2$. This capacity exclusion is independent
of the earlier source/initial-run gap exclusions.

\begin{theorem}[The entire both-nonincreasing first half]
For two actual nonrectangular increasing arms, suppose BOTH forward
factors are nonincreasing at $j+1$ and
\[
                             2j\le a+b+2.           \tag{8}
\]
Then $T\le0$. No source-prefix, nonnegative-tail, OLD A2,
head, or product-mode assumption is needed.
\end{theorem}
\begin{proof}
Orient $a\le b$. The preceding theorem handles $j\le a$.
If $a+2\le j\le b$, then $r=a-j\le-2$, so $C_Z=0$,
and the first-height proposition gives $T\le0$.
If $j=a+1\le b$, then $C_Z=1$. The opposite nonincreasing
premise forces $\kappa_Y\le j+1=a+2$: otherwise both opposite
coefficients are still universal and strictly double.
Thus $0\le\kappa_Y-2\le a$, and the first-height proposition
pays $C_Z$ by $q_Z(\kappa_Y-2)\ge1$.

It remains that $j>b$. Integrality, (8), and $a\le b$ force
$a=b$, $j=a+1$, $n=a-1$. All source-layer products still
have nonnegative drops, because their endpoint sums are at most
$a+b-1=2n+1$. The lowest source height interval contributes
\[
 q_Z(\kappa_Y-2)-q_Z(a)=q_Z(\kappa_Y-2)-1.
\]
If $g$ is the opposite non-one count, then
$3\le\kappa_Y\le g+2$ and $2g\le a-1$ imply
$1\le\kappa_Y-2\le a-2$.
Every interior coefficient of $q_Z$ on $[1,a-2]$ is at
least two: $q_Z(1)=h+1\ge2$, $q_Z(a-2)=A_2=2$, and
unimodality bounds every intervening coefficient below by
the smaller endpoint value. The displayed drop is therefore
at least one and pays $C_Z=1$.
\end{proof}

The nonrectangular hypothesis in this first-half theorem is real.
For two equal all-ones arms of weight $a$, the triangular product
has $T=1$ at $j=a+1$, which satisfies (8); that edge is after
the ordinary product mode and therefore is not a proper S edge.

\section{The reverse half of a product tail}

\noindent{\small Source\enspace\path{agent_half_lc_actual_weak_curvature.tex}.}\par\medskip

For a polynomial $p$, zero-extend its coefficients and write
\[
 E_p(n)=2p_n-p_{n-1},\qquad
 T_p(n)=E_p(n+1)-E_p(n)=2p_{n+1}-3p_n+p_{n-1}.
\]
\begin{theorem}[Ordinary half-tail source nonrecovery]\label{thm:actual-half-tail}
Let $X,Y$ be normalized actual increasing nonempty arms, and put
$p=q_Xq_Y$, $D=\operatorname{wt}X+\operatorname{wt}Y$.
If $\nu$ is the LAST ordinary coefficient mode of $p$, then
\[
                        T_p(n)\le0\qquad(\nu<n,\ 2n\le D). \tag{W}
\]
No nonnegative-source hypothesis is required. In fact $E_p(n)\ge0$
through this whole half range independently of (W).
\end{theorem}

\paragraph{Actual inputs used in the proof.}
A normalized arm $Z=1^\alpha U$ has non-one count $h$, weight $w$,
ordinary polynomial $q_Z$, main interval product
$K_Z=[\alpha+1]\prod_{u\in U}[u]$ of degree $w-h$, and corrected source
\[
 e_Z=(2-z)q_Z-z^w+z^{w+1}.
\]
The established actual inputs are:
\begin{itemize}
\item $q_Z$ is positive on $[0,\operatorname{wt}Z]$, unimodal,
and every ordinary height interval has endpoint sum at most its weight;
\item $e_Z$ is nonnegative unimodal, positive exactly on
$[0,\operatorname{wt}Z+1-\kappa_Z]$, starts with coefficient two,
and every height interval has endpoint sum at most $\operatorname{wt}Z-1$;
\item $q_Z\succeq K_Z$ and
$(q_Z(j)-q_Z(j-1))^+\le K_Z(j)$ at every rank;
\item $2h\le\operatorname{wt}Z-1$ and
$\operatorname{wt}Z\ge\kappa_Z(\kappa_Z-1)/2$;
\item for a nonrectangular arm, its second ordinary height interval is
exactly $[1,\operatorname{wt}Z-2]$.
\end{itemize}
The main-box derivative bound is
{\normalfont\path{main_box_derivative_envelope.tex}}; the mass bound is
{\normalfont\path{first_defect_mass.tex}}. The remaining shape inputs are the
established ordinary and corrected-source layer theorems.

\paragraph{A window identity and its crossing case.}
Suppose $p$ is unimodal, with first mode $m$ and last mode $\ell$,
and $w=[t]p$. On the ordinary tail of $w$, put $k=n+1-t$.
Then
\[
 T_w(n)=E_p(n+1)-E_p(k),\qquad
 \Delta w_{n+1}-\Delta w_n=\Delta p_{n+1}-\Delta p_k.   \tag{1}
\]
The first window mode is at least $m$. If $k\le\ell$, the leaving
input increment is nonnegative, and the entering input increment is
nonpositive because $n$ is after the window mode and therefore at least
the first old mode. Hence $\Delta w_{n+1}\le\Delta w_n$; since the window
is on its ordinary tail, $\Delta w_{n+1}\le0$ as well. Therefore
\[
 T_w(n)=\Delta w_{n+1}+(\Delta w_{n+1}-\Delta w_n)\le0. \tag{2}
\]
This includes all maximum plateaus and negative leaving indices.
Only the case $k>\ell$ needs a source comparison.

\begin{lemma}[Rectangle-factor base]
Statement (W) holds whenever at least one arm is rectangular.
\end{lemma}
\begin{proof}
Write $p=[A+1]q$ with $\deg q=b$.
If $A\ge b$, its maximum plateau is exactly $[b,A]$, so there is no
strict post-LAST-mode position through $(A+b)/2$.
If $A<b$, write $e=(2-z)q-z^b+z^{b+1}$. Its proved unimodality and
the causal averaging relation $2q_j=e_j+q_{j-1}$ below rank $b$ imply
that the source mode is no later than the first $q$ mode. Consequently
the raw source of $q$ is nonincreasing after its ordinary mode up to
rank $b-1$. The window argument (1)--(2) pays every required comparison
with $n+1<b$: in its noncrossing case both source endpoints lie in
this monotone prefix.

The only remaining possibility is $n+1=b$, which the half-range
condition permits only for $b-A\in\{1,2\}$. The exact identity
\[
                  (2-z)p=[A+1]e+z^b-z^{A+b+1}
\]
gives
$T_p(b-1)=1-e_{b-A-1}\le0$, since $e_0=2$ and $e_1\ge1$.
Two rectangular factors can also be handled directly by their
log-concave trapezoidal product.
\end{proof}

\paragraph{Induction setup and the NEW terminal-source edges.}
Assume both new arms are nonrectangular. Choose the NEW heavier arm,
of weight $a\ge b$, delete its last non-one part $t$, and call its
old weight $c=a-t$. Let $p$ be the old ordinary product, and $P$ the
new one. Then
\[
                          P=[t]p+z^a q_Y.              \tag{3}
\]
The total non-one count decreases by one, so (W) holds for $p$ by
induction. The first mode of $P$ is strictly below $a$.
At that mode $P$ agrees with $w=[t]p$, and $w\le P$ everywhere,
so $\max w=\max P$ and every maximum of $w$ is also a maximum of $P$.
In particular
$\operatorname{lastmode}w\le\operatorname{lastmode}P$.
Thus a strict ordinary-tail position of $P$ is on the window tail.

The source in (3) can affect a half-range coefficient triple only when
$n\ge a-1$, hence only for NEW weight gaps $a-b=0,1,2$.
Those edges are all covered by the proved theorem ``the entire
both-nonincreasing first half'' in
{\normalfont\path{agent_bothdeclining_first_birth.tex}}:
\[
\begin{array}{c|c|c}
a-b&n&j=a+b-n\\\hline
2&a-1&a-1=b+1\\
1&a-1&a\\
0&a-1,\ a&a+1,\ a
\end{array}
\]
At each listed forward edge both actual factors are nonincreasing,
and $2j\le a+b+2$. That theorem gives $T_P(n)\le0$ without a mode,
source-sign or OLD assumption. Henceforth
\[
                              n+1<a.                   \tag{4}
\]
Then $T_P(n)=T_w(n)$ and $k=n+1-t<c$.
By (2) a possible failure must have $k>\ell$, the LAST old mode.

\paragraph{An OLD rectangular arm created by deletion.}
The unchanged old opposite arm is nonrectangular. If the deleted arm
has become the rectangle $[c+1]$, the old product is $[c+1]q_Y$ and
its LAST mode is at least $c$: it strictly increases through rank $c$
when $c<b$, and has maximum plateau $[b,c]$ when $c\ge b$.
But (4) gives $k<c\le\ell$, so (2) pays the edge.
The remaining induction case therefore has BOTH old arms nonrectangular.

\paragraph{The OLD positive source prefix and the entering correction.}
Orient the OLD weights as $M\ge s$, write $d=M+s$, and let
$\kappa\ge3$ be the larger old arm's first missing weight. Put
\[
 R=M+1-\kappa,\quad
 C=e_Lq_S,\quad E=(2-z)p.
\]
Then
\[
                         E=C+z^M(1-z)q_S.              \tag{5}
\]
Thus $E_j=C_j\ge0$ below $M$.
Every interval-product component of $C$ has endpoint sum at most
$M+s-1=d-1$, so $C$ is nonincreasing from $\lfloor d/2\rfloor$ onwards.
Together with the old induction hypothesis, this proves
\[
 E_u\ge E_v\quad(\ell<u\le v<M).                       \tag{6}
\]
No global sign-prefix or OLD A1 assertion is used.

If $n\le M-2$, both source endpoints in (1) lie in this old prefix,
so (6) pays the edge. Suppose now $n\ge M-1$.

\begin{lemma}[The initial entering-source boundary]
The position $n=M-1$ is always paid.
\end{lemma}
\begin{proof}
Here the entering value is $E_M=C_M+1$.
If $\kappa\le s+2$ and $M>s$, pair the lowest source layer $[0,R]$
with the lowest ordinary layer $[0,s]$. Its trapezoid loses exactly
one at $M-1\to M$, because the leaving point $\kappa-2$ lies in
$[0,s]$. All other component differences are nonnegative, so
$E_M\le E_{M-1}\le E_k$.
If $M=s$, use the second ordinary layer $[1,s-2]$ instead: the leaving
point $\kappa-2$ belongs to it because $3\le\kappa\le M$.

If $\kappa>s+2$, then $C$ has support endpoint
$R+s=M+s+1-\kappa\le M-2$, so $E_M=1$.
The half-range condition gives $k=M-t\le s+2$.
Since $s\ge3$, this case has $\kappa\ge6$, and first-missing mass
gives $R\ge2$. The convolution $C$ is positive on all of $[0,R+s]$,
which contains $[0,s+2]$. Since $k>\ell\ge0$ and $k<M$,
$E_k=C_k\ge1=E_M$.
\end{proof}

\paragraph{All later entering-source positions.}
Write $n=M+r$, $r\ge0$. Equations (1) and (5) give
\[
 T_P(n)=C_{M+r+1}+q_S(r+1)-q_S(r)-E_k.                 \tag{7}
\]
The new half-range condition gives
\[
                      k\le k_*:=s-r+1.                 \tag{8}
\]
Every comparison of a component at $M-1$ with $M+r+1$ has nonnegative
drop, since the two indices sum to $2M+r\ge d-1$.

\begin{lemma}[Late entering-source payment]
If $r+3\ge\kappa$, then (7) is nonpositive.
\end{lemma}
\begin{proof}
For each $q_S$ height interval $[A,B]$ starting at $A=r+1$, pair it
with the lowest source layer $[0,R]$. At $M-1$ this convolution counts
the entire interval, because $B\le s-r-1\le M-1$ and
$M-1-R=\kappa-2\le A$. At $M+r+1$ it misses at least $A$, since
$M+r+1-R=r+\kappa\ge A+1$.
Thus the total drop is at least the number of such starting layers,
which is at least $q_S(r+1)-q_S(r)$. Therefore
$E_{M+r+1}\le E_{M-1}\le E_k$, by (6).
\end{proof}

Only $0\le r\le\kappa-4$ remains. Let $h_S$ be the old short arm's
non-one count and $m_S=s-h_S$ the degree of its main interval product.

\begin{lemma}[A short main-box payment]
If $m_S\ge\kappa+r-1$, then (7) is nonpositive.
\end{lemma}
\begin{proof}
The lowest source layer contributes exactly
\[
 \sum_{u=\kappa-2}^{\kappa+r-1}q_S(u)
                  -\sum_{u=M}^{M+r+1}q_S(u)             \tag{9}
\]
to $C_{M-1}-C_{M+r+1}$. The second sum is zero if $M>s$, and one if
$M=s$. Let $K_S$ be the short arm's main interval product. Its symmetry
and unimodality, together with $r+1\le\kappa-2$ and
$\kappa+r-1\le m_S$, give
\[
 q_S(\kappa-2)\ge K_S(\kappa-2)\ge K_S(r+1)
                         \ge(q_S(r+1)-q_S(r))^+.
\]
There is at least one further positive coefficient in (9), namely
$q_S(\kappa-1)\ge1$, since $\kappa-1\le m_S\le s$.
It pays the possible terminal penalty one. Other component drops are
nonnegative, so again $E_{M+r+1}\le E_{M-1}\le E_k$.
\end{proof}

\begin{lemma}[A virtual central head payment]
Suppose $k_*<M$ and $\kappa\le M-s+2r+1$. Then (7) is nonpositive.
The number $M-s+2r$ is only a comparison width, not a claimed legal arm part.
\end{lemma}
\begin{proof}
Compare $C_{k_*}$ with $C_{M+r+1}$. Their indices sum to $d+2$,
so every component has nonnegative drop. For a $q_S$ height interval
beginning at $A=r+1$, its right endpoint obeys
$B\le s-r-1=k_*-2$. At the head $k_*$ the lowest source layer counts
all of $[A,B]$, since
\[
                   k_*-R=s-r+1-(M+1-\kappa)\le A.
\]
At the tail it misses $A$ as in the late payment.
Thus $E_{M+r+1}\le C_{k_*}=E_{k_*}\le E_k$ by (6) and (8).
\end{proof}

\paragraph{The two resources cover every nonexceptional head.}
Suppose $k_*<M$. If both preceding payments failed, then
\[
 M\le s+\kappa-2r-2,\qquad m_S\le\kappa+r-2.
\]
Normalization $2m_S\ge s+1$ gives $s\le2\kappa+2r-5$, and hence
$M\le3\kappa-7$. This contradicts first-missing mass:
\[
 \kappa(\kappa-1)/2>3\kappa-7
       \quad(\kappa\ge3),
\]
because the doubled gap is
$(\kappa-3)(\kappa-4)+2>0$ for every integer $\kappa\ge3$.

\paragraph{The three exceptional heads are also covered.}
The condition $k_*=s-r+1\ge M\ge s$ has only three possibilities.
\begin{itemize}
\item $M=s,r=0$: mass gives $s\ge2\kappa-3$, so
$2m_S\ge s+1$ implies $m_S\ge\kappa-1$ and the main-box payment applies.
\item $M=s+1,r=0$: mass gives $M\ge2\kappa-3$, and
$2m_S\ge s+1=M$ implies the same integer bound $m_S\ge\kappa-1$.
\item $M=s,r=1$: for $\kappa\le4$ the late payment already applies;
for $\kappa\ge5$, mass gives $s\ge2\kappa-1$, hence
$m_S\ge\kappa$ and the main-box payment applies.
\end{itemize}
This pays every source-free half-range comparison after the last new mode.
Together with the NEW terminal-source edges and the rectangular cases,
the induction is complete, proving (W).


\section{The closed interval comparison}

\noindent{\small Source\enspace\path{coupled_source_window_A2_and_kappa_nine_original_proof.tex}.}\par\medskip

\begin{theorem}[A complete source-to-window A2 bridge]
Let $p$ be an actual old ordinary arm product. Assume FULL OLD A2
and coefficient tilt A1 for $p$. Put $u=(\Delta p)^+$, let $m$ be
the LAST ordinary mode of $p$, and let $\ell$ be the first mode of
$2^nu_n$. Set $E_n=2p_n-p_{n-1}$.
Assume the following CLOSED-gap proper comparison for every $t\ge1$:
whenever $n>m$, $k=n-t\ge0$, $k<\ell$,
\[
 p_n>p_k,\quad p_{n+1}>p_{k+1},\quad 0\le E_n\le E_k,
\]
one has $E_{n+1}-E_n\le E_{k+1}-E_k$.
Then EVERY ordinary uniform window $[t]p$ has FULL A2.
This theorem concludes weighted-positive-D UI, not whole raw-source UI.
\end{theorem}
\begin{proof}
Put $v=[t]p$ and $w_n=\Delta v_n=p_n-p_{n-t}$. Uniform convolution
makes $v$ ordinarily UI. Its positive differences have initial interval
support: before the old first mode the entering term strictly exceeds
the outgoing term; across the old mode interval its differences are
nonincreasing; after that interval they are nonpositive. This also
follows directly from the established actual positive-difference support.
Write $\sigma$ for its first mode and $f_n=2^n(w_n)^+$.

First build a COMPLETE UI head comparison, including its first boundary.
For $n\le m$, every signed old difference is nonnegative, so
$w_n=([t]u)_n$. Its weighted sequence is convolution of the old A2
sequence by $(1,2,\ldots,2^{t-1})$, a log-concave kernel without holes.
Such convolution preserves UI. Denote this WHOLE UI sequence by $g$.
At $m+1$, signed convolution and clipping give $f_{m+1}\le g_{m+1}$.
Alternatively, for $n<t$ one has $w_n=p_n$ and $f_n=2^np_n$,
whose whole sequence is UI by A1; at $t$, $f_t\le2^tp_t$.
Choose
\[
                         c=\max(m,t-1).
\]
If $m\ge t-1$ use the first head sequence; otherwise use the second.
In either case $f=g$ through $c$, $f_{c+1}\le g_{c+1}$, and $g$ is
WHOLE UI. An earlier strict fall of $f$ therefore forces
$f_{c+1}\le f_c$. It remains to propagate NONSTRICT falls after this
boundary; using strict falls alone would leave an equality gap.

At a crossing $n>m$, $n\ge t$, if $w_n,w_{n+1}>0$, set $k=n-t$.
If $k\ge\ell$, the old weighted geometric tail gives
\[
 w_n\le\sum_{i=k+1}^{\mathrm{firstmode}(p)}u_i\le2u_{k+1}.
\]
The finite sum bound follows from $2u_{i+1}\le u_i$ throughout the
used tail; its sharper telescoping form subtracts the last coefficient.
Also
\[
 w_{n+1}=w_n-u_{k+1}-(p_n-p_{n+1}),\qquad p_n-p_{n+1}\ge0,
\]
so $2w_{n+1}\le w_n$ unconditionally in this branch.

If $k<\ell$, then $E_k,E_{k+1}>0$ and
$H=E_{k+1}-E_k=2u_{k+1}-u_k\ge0$. A nonstrict weighted fall means
\[
          2w_n-w_{n-1}=E_n-E_k\le0.                    \tag{1}
\]
If $E_n<0$, A1 implies $E_{n+1}\le0$, hence
$2w_{n+1}-w_n=E_{n+1}-E_{k+1}<0$.
Otherwise the CLOSED-gap hypothesis and (1) give
\[
 E_{n+1}-E_{k+1}\le E_n-E_k\le0,
\]
again proving $2w_{n+1}\le w_n$. Once the positive-difference support
ends, $f$ is zero and every later edge is nonincreasing.
Thus nonstrict falls propagate throughout the entire crossing range.
The kernel-checked head/boundary/closed-cross lemma excludes every
strict fall followed by a later rise of $f$. Finiteness supplies a
whole up/down mode. This proves FULL A2 for $[t]p$.
\end{proof}

\paragraph{The relative raw-kernel theorem retains a strict surplus at equality.}
The frozen relative-head proof assumes $0\le E_n<E_k$, but its exact
lower bound is stronger: with $\lambda=C/d_r>0$,
\[
 H\ge\lambda E_k+1,\qquad E_n\ge d_r>0,\quad T\le C.
\]
Thus $E_n\le E_k$ already gives $H\ge C+1>T$. Its proof of the
raw-kernel convolution inequality used neither strictness of the gap
nor C. The CLOSED-gap extension is therefore valid, with no appeal to
an unproved continuity or equality argument. It covers every $r\ge k$.

\section{Small first missing weights}

\noindent{\small Source\enspace\path{coupled_source_window_A2_and_kappa_nine_original_proof.tex}.}\par\medskip

\begin{theorem}
Every final-deletion-closed class of ACTUAL normalized increasing arms
with FULL C and $\kappa\le9$ has FULL pair A2. The ORIGINAL strict
core and BOTH relatives are coefficient UI at EVERY strict peak,
including ONE arbitrary legal heavier append.
For an old nonrectangular pair with maximum first missing weight nine,
EVERY ordinary uniform window of its reverse product has FULL A2.
No whole raw-source-window UI conclusion is asserted.
\end{theorem}

\subsection{All positive mixed recoveries have shared forward rank at most ten}
The old maximum-$\kappa\le8$ class and arbitrary heavier append are
frozen. Consider an old pair of maximum first missing weight nine,
with FULL OLD A2 in the count induction. A1 follows from FULL C.
Both-rising C and both-nonrising actual resources give $T\le0$ after
the LAST ordinary product mode whenever the current source is nonnegative.
The latter includes both halves by combining the unconditional forward
both-nonincreasing theorem and the actual reverse half-tail theorem.
Only mixed recovery remains.

If its rising opposite has $\kappa\le8$, the complete frozen pointwise
exceptions have owner weight at most eight. An owner of first missing
weight nine has weight at least 36, so none is possible. Hence a positive
new mixed recovery must have opposite first missing weight NINE.
With $Z=1^\alpha U$ the nonincreasing owner, use the exact initial
exchange $\alpha_* =\max(\alpha,t_1-1)$.
The frozen absolute-nine theorem pays all owners with $h\ge6$.
The gap-six theorem pays $h=4,5$ when $\alpha_*\ge3$; the gap-five
theorem pays $h=2,3$ when $\alpha_*\ge4$.
The remaining canonical domain is
\[
\begin{array}{c|c|c|c}
 \alpha& h& t_1& D_{\max}\\\hline
 1&2,3,4,5&2&71\\
 2&2,3,4,5&2,3&56\\
 3&2,3&2,3,4&36
\end{array}                                                     \tag{2}
\]
Here $D=\sum(t_i-1)$, $d=9-\alpha$, and the unbounded necessary
bound is $D\le\min(2d(d-3),12d-25)$.

For $r=a-j\ge0$, positive recovery forces positive primitive curvature,
$D=2r+2+e$, and $0\le e\le d-3$. Put $L=d-e-3$ and
\[
 R_m=\frac{\binom{r+e+m+2}{m}}{\binom r m}.
\]
The genuine compulsory source increments give the frozen sufficient
rank payment
\[
 \frac{\mathbf1_{L\le r}}{R_L}
 +\sum_{\substack{2\le v\le6\\L+v\le r}}
       \frac{2^{v-2}}{R_{L+v}}\ \ge\ \frac{e+1}{r+1}.    \tag{3}
\]
All these offsets are in the mode-free early-source range through
$\kappa-3=6$; included birth reads have nonnegative actual indices.
The old CI-relative/reflection proof of (3) applies for every $h\ge2$,
not only the $h\ge6$ specialization of the earlier nine-threshold result.

The COMPLETE rank grammar (2)--(3) has 406 states. Exact rational
arithmetic pays 291. Its remaining 115 states have $D\le25$.
For EVERY remaining state enumerate all sorted partitions of $D$ into
the specified counts, with minimum increment at most $\alpha$.
This reconstructs 2143 distinct canonical owners and ALL 4377 used
owner/rank inputs. Of these, 3757 have nonpositive curvature and are
irrelevant to positive recovery. At the 620 positive-curvature inputs,
the direct complete compulsory-source bound pays 601:
\[
 F_Z(v)+F_Z(v-2)+2F_Z(v-3)+4F_Z(v-4)
                  +8F_Z(v-5)+16F_Z(v-6)\ge C,
 \qquad v=a-r-8.                                      \tag{4}
\]
The ONLY 19 source-unpaid candidates all have $j=a-r\le10$.
All rational numerators/positive denominators, nonpositive-curvature
classifications, payments (4), and candidate rank bounds reduce to
719 distinct nonnegative integer obligations in six ordinary-kernel
chunks and ONE complete public aggregate. An independent degree/maximum-part
grammar must reconstruct the entire rank and owner partition, every
category, every numerator, and the public obligations.

For $r\le-2$ no primitive curvature is available. At $r=-1$, the
unit birth is absent only when $a+2-9<0$, so $a\le6$ and $j=a+1\le7$.
One-part owners have positive primitive curvature ONLY at this same
negative index and are handled by the identical bound. Therefore
\[
             T>0\quad\Longrightarrow\quad j\le10       \tag{5}
\]
throughout the entire new mixed/source-prefix region, independently
of every opposite capacity, count, padding and weight.

\subsection{A genuine strict opposite head pays all nineteen candidates}
At a proper crossing, BOTH strict positive windows give
\[
                  k+n+2<a+b,\qquad k+1\le j-2\le8.
\]
This uses actual upper-centre reflection of $q_Zq_Y$, not a synthetic
scalar rank profile. Also $j\le a+1$, hence $k+1\le a-1$.
The entire used RAW owner source $R_Z=(2-z)q_Z$ is consequently
nonnegative, because it agrees with the corrected positive source
below degree $a$; its constant coefficient is two.

The opposite main interval product has degree $M$ and largest factor
degree $B$. Its complete first-missing staircase forces
\[
 M\ge1+\frac{7\cdot8}{2}=29,\qquad
 M-B\ge1+\frac{6\cdot7}{2}=22.
\]
The exact CI positive-primitive support formula gives positive
differences through $\min(\lfloor M/2\rfloor,M-B)\ge14$.
Every other actual shifted-CI summand has endpoint sum at least $M$,
so its differences are nonnegative through $\lfloor M/2\rfloor$.
Thus the ACTUAL opposite coefficients satisfy
\[
                       \Delta q_Y(i)\ge1\quad(0\le i\le14).      \tag{6}
\]
No coefficientwise lower bound is differentiated to obtain (6).

The exact head increment is $H=[z^{k+1}]R_Z(1-z)q_Y$.
All used terms are nonnegative. If $0\le r<k$, its owner term $r+1$
and (6) give
\[
 H\ge R_Z(r+1)\Delta q_Y(k-r)
   \ge R_Z(r+1)=q_Z(r+1)+d_{r+1}\ge C\ge T.
\]
For $r=-1$ the owner-constant term gives $H\ge2\ge T$.
For $r\ge k$, the CLOSED-gap relative raw-kernel proof above gives
$H>T$. Every positive recovery is therefore paid at EVERY proper
closed-gap head. For $T\le0$, FULL OLD A2 gives $H\ge0$.
Terminal forward ranks zero/one are paid directly or have negative
current source. This proves the complete closed-gap source criterion.

\subsection{FULL induction and original coefficients}
The first theorem gives FULL A2 for EVERY uniform window of an old
maximum-nine pair. In the non-one-count induction orient a heavier
new arm $Zt$. Empty, rectangular and heavier-one-part bases are frozen.
If the old pair has maximum first missing weight at most eight, its
frozen arbitrary heavier-append theorem applies. Otherwise use the new
window theorem. The exact actual correction starts at the new heavier
weight, and the NEW first ordinary mode is strictly below that weight.
Hence its positive differences are a terminal truncation of the window's
positive differences, followed by zeros; FULL A2 is preserved.
This closes the induction and also proves ONE arbitrary extra legal
heavier append without requiring C of that new arm. The frozen proper
deletion and lower-peak/reflection joins then give ORIGINAL coefficient
UI of the strict core and BOTH relatives.

Let A15 be A14 together with ALL max-four-prefix arms followed by FOUR
successive dominant capacities $T_i\ge2(W_0+\sum_{j<i}T_j)$, with
$T_1\ge\max(5,2W_0)$, and all max-eight tails of nonbinary weight at
most 32. Binary counts and initial padding are arbitrary. Frozen
dominant C and max-four/fixed-weight C give FULL C; final deletion
closure and $\kappa\le9$ are exact. Thus BOTH original arms may
independently have FOUR such jumps, arbitrary strict peak, and ONE
further legal heavier append. No amplification is used.


\section{First missing weights ten through twelve}

\noindent{\small Source\enspace\path{first_missing_twelve_original_criterion_proof.tex}.}\par\medskip

All notation and actual-model lemmas are those of the frozen
{\normalfont\path{coupled_source_window_A2_and_kappa_nine_original_proof.tex}}.
In particular FULL C, FULL A1 and FULL A2 have their original coefficient
meanings. No conclusion about raw-source-window unimodality is used here.

\begin{theorem}
Every final-deletion-closed class of normalized ACTUAL increasing arms
with FULL C and first missing weight at most twelve has FULL pair A2.
Its ORIGINAL unamplified strict core and BOTH relatives have unimodal
coefficients at EVERY strict peak. ONE further arbitrary legal append
to the heavier arm is also permitted without assuming C of that new arm.
\end{theorem}

\subsection{The unbounded reduction and the complete finite domain}
Work in the count induction, assuming FULL OLD A2 for an old pair.
Let $p=q_Zq_Y$, $a=\operatorname{wt}Z$, $b=\operatorname{wt}Y$,
$j=a+b-n$, $r=a-j$, $E_i=2p_i-p_{i-1}$, and
$T=E_{n+1}-E_n$. The relevant proper CLOSED-gap comparison is
\[
 n>\operatorname{lastmode}(p),\quad k<\operatorname{firstmode}(2^i(\Delta p_i)^+),
 \quad p_n>p_k,\quad p_{n+1}>p_{k+1},\quad 0\le E_n\le E_k.
\]
Both-forward-rising C and both-forward-nonincreasing actual theorems
cover their entire respective phases. In the mixed phase orient $Z$ as
the nonincreasing arm. The opposite source prefix is nondecreasing at
all used ranks by the frozen slope/source mode ordering. Write
$K=\kappa_Y$. A nonpositive $T$ is paid by the nonnegative OLD-A2
head increment, so consider $T>0$.

The frozen nine-threshold proof handles $K\le9$ when both first
missing weights are at most nine. If instead the owner has first missing
weight at least ten, it has at least eight non-one parts. The absolute
eight theorem for four-part owners and absolute nine theorem for six-part
owners therefore exclude positive recovery for EVERY $K\le9$ in this
new situation. It suffices to treat $10\le K\le12$ and arbitrary owners.

Use the exact initial exchange to impose $t_1-1\le\alpha$ for
$Z=1^\alpha(t_1,\ldots,t_h)$ without changing its weight, count or
polynomial. For $r\ge0$, positive recovery requires
\[
 \begin{gathered}
 d=K-\alpha\ge6,\qquad D=\sum_i(t_i-1)=2r+2+e,\\
 0\le e\le d-3,\qquad D\le D_{\max}:=\min(2d(d-3),12d-25).
 \end{gathered}
                                                        \tag{1}
\]
The gap-six theorem excludes $h\ge4$ at $d\le6$.
The UNBOUNDED gap-eight theorem for six-part owners excludes
$h\ge6$ at $d\le8$, at every $K$, not just $K=9$. Thus
\[
 1\le\alpha\le K-6,\qquad
 2\le h\le h_{\max}:=\min\left(D,
 \begin{cases}3&d\le6,\\5&7\le d\le8,\\D&d\ge9.\end{cases}\right).
                                                        \tag{2}
\]
All parameters discarded by (1)--(2) are covered by the frozen
UNBOUNDED lemmas; no capacity or count cutoff is a premise.

Put
\[
 M_K=1+(K-2)(K-1)/2,\quad O_K=1+(K-3)(K-2)/2,
 \quad S_K=\lfloor M_K/2\rfloor.
\]
Their exact values are
\[
\begin{array}{c|rrr}K&M_K&O_K&S_K\\\hline
10&37&29&18\\11&46&37&23\\12&56&46&28
\end{array}.                                             \tag{3}
\]
For a rank state (1), let $j_{\max}=\alpha+D+h_{\max}-r$.
If $j_{\max}-2\le S_K$, every owner in that state has its proper
head inside the strict opposite prefix established below. Such a state
needs no source payment.

For every remaining state apply the genuine source-sum/CI-ratio test.
Set $L=d-e-3$ and
\[
 b_m=\begin{cases}
 \binom r m/\binom{r+e+m+2}m&0\le m\le r,\\0&\text{otherwise}.
 \end{cases}
\]
The following sufficient condition pays the entire positive curvature
$C=(\Delta q_Z)_{r+1}-(\Delta q_Z)_r$:
\[
             b_L+\sum_{v=2}^{K-3}2^{v-2}b_{L+v}
                           \ge\frac{e+1}{r+1}.          \tag{4}
\]
It is the frozen arbitrary-count source-ratio argument, with ALL
available early mode-free source increments. It does not differentiate
a lower bound for source coefficients.

The COMPLETE rank grammar (1)--(4) gives the following exact partition:
\[
\begin{array}{c|rrrr}
 K&\text{all ranks}&\text{strict head}&\text{source (4)}&\text{owner complement}\\\hline
10&684&170&451&63\\
11&1048&279&641&128\\
12&1509&410&891&208\\\hline
\text{total}&3241&859&1983&399
\end{array}.                                             \tag{5}
\]
For each complementary rank, any owner not already inside the strict
head must have
\[
 h\ge h_{\min}:=\max(2,S_K+3-\alpha-D+r).              \tag{6}
\]
Enumerate ALL partitions of $D$ into $h_{\min}\le h\le h_{\max}$
positive sorted increments, with minimum at most $\alpha$.
The largest remaining $D$ is respectively $33,39,47$ for $K=10,11,12$.
This is the complete complement of the proved unbounded reductions,
including every owner rank, rather than a chosen sample of arms.

Construct each actual $F_Z$ and compute the signed primitive and curvature.
There are 1,858,242 owner/rank inputs, including 945,847 distinct owners
when counted separately at each $K$. Exactly 1,786,199 have negative
primitive or nonpositive curvature. ALL 72,043 positive-curvature inputs
are paid by the actual source sum
\[
 C\le F_Z(v)+\sum_{i=2}^{K-3}2^{i-2}F_Z(v-i),
                  \qquad v=j+1-K,                       \tag{7}
\]
using zero extension. By $K$ the positive-curvature counts are
$4877,17998,49168$. There is NO unpaid owner outside the strict-head range.

The independent verifier enumerates ranks by total degree and reverse
rank, and partitions by MAXIMUM part, whereas the builder enumerates
source gap and excess, and partitions by MINIMUM part. It reconstructs
the actual polynomial as a sum of positive suffix boxes, with a closed
binomial formula for its binary block; it does not reuse the builder's
chronological actual-arm recurrence. Product fractions independently
reconstruct the rational ratios in (4). It matches every noncritical
classification, every critical coefficient operand, all 399 layer hashes,
the complete state hash, and every line of the 72,043-row critical stream.

The explicit Lean certificate contains ALL 72,043 source SUM inequalities
(7), 3266 distinct integer numerator/positive-denominator tests for (4),
and 63 distinct strict-head rank pairs. These are 75,372 obligations in
754 chunks, compiled in 31 memory-bounded modules, and ONE final public
aggregate containing EVERY module's complete conjunction. Noncritical curvature classifications
and exhaustion of the word grammar are independently exact Python checks,
not an end-to-end Lean implementation of words or their enumerator.
Four further generic Lean reports check (2), (3), proper-head safety,
and all negative-birth guards. Standard axioms only are required.

At $r\le-2$ positive curvature is absent. At $r=-1$ a source-unpaid
unit birth requires $a+2<K$, so $j=a+1\le K-2$. This includes the
only possible one-part-owner birth exception. By (3),
$j-2\le K-4\le S_K$. Consequently the complete conclusion is
\[
 T>0\quad\Longrightarrow\quad j-2\le S_K
          \qquad(10\le K\le12).                       \tag{8}
\]

\subsection{Strict actual heads, closed-gap compensation and the full induction}
The BOTH positive windows imply $k+n+2<a+b$, hence
$k+1\le j-2\le S_K$. Also $j\le a+1$, so $k+1\le a-1$.
The used raw owner source $R_Z=(2-z)q_Z$ therefore agrees with the
nonnegative corrected source throughout the entire used prefix.

The exact opposite first-missing staircase gives its main interval product
degree $M\ge M_K$, and degree outside its largest factor at least $O_K$.
The CI positive-primitive support formula and $2S_K\le M_K$,
$S_K\le O_K$ give strictly positive main-box differences through $S_K$.
EVERY proper shifted-CI source has endpoint sum at least $M$, so its
differences are nonnegative through that head, including birth boundaries.
Thus the ACTUAL, SIGNED opposite differences satisfy
\[
                       \Delta q_Y(s)\ge1\quad(0\le s\le S_K).   \tag{9}
\]
This uses genuine CI differences, not differentiation of a coefficient bound.

For $H=E_{k+1}-E_k=[z^{k+1}]R_Z(1-z)q_Y$, all used summands
are nonnegative. If $0\le r<k$, select the owner index $r+1$:
\[
 H\ge R_Z(r+1)\Delta q_Y(k-r)
 \ge q_Z(r+1)+(\Delta q_Z)_{r+1}\ge C\ge T.
\]
At $r=-1$ select the constant owner term to obtain $H\ge2\ge T$.
For $r\ge k$, the frozen raw-relative LC kernel retains its strict
$+1$ surplus and applies to CLOSED $E_n\le E_k$, including equality.
All positive recoveries are paid. Nonpositive recoveries use $H\ge0$
from FULL OLD A2. This proves the entire closed-gap proper S criterion.

The frozen WHOLE closed-gap S-to-window theorem now proves FULL A2 for
EVERY $[t]p$. In the count induction the old maximum-nine case is frozen;
the new maximum-ten-through-twelve case uses the criterion just proved.
Orient the newly appended arm as the heavier arm. Its actual terminal
correction starts at its new weight, strictly after the new first ordinary
mode. Hence positive differences are a truncation of the uniform window's
positive differences. FULL A2 follows, as does ONE arbitrary further
legal heavier append. Frozen proper deletion, lower peak, reflection and
relative joins give ORIGINAL coefficient unimodality at every strict peak.


\section{Dyadic bounds and a strict earlier increment}

\noindent{\small Source\enspace\path{actual_dyadic_slope_thinness_and_high_source_mixed_S_proof.tex}.}\par\medskip

Set $\rho=3/5$, $b=35/32$, $c=3/2$, $q=b^2/(2\rho)=6125/6144$,
and $v=b/(2\rho)=175/192$. For $t\ge2$ put
\[
 s_t=S_t(\rho)=\tfrac32(1-\rho^t).
\]
All capacities, counts, weights and initial runs below are UNBOUNDED.

\subsection{A uniform positive evaluation bound, including the whole append fold}
For EVERY formal word $1^\beta V$, allowing $\beta=0$, let
$M=\beta+\sum_{t\in V}(t-1)$. Then
\[
                       F_{1^\beta V}(\rho)\le10b^M-125/22<10b^M. \tag{1}
\]
The bound does not require ordered capacities. To prove it, first note
\[
 0\le s_t\le b^{t-1}\quad(t\ge2),\qquad
 s_t\ge147/125=1+22/125\quad(t\ge3).                  \tag{2}
\]
For $t=2,3,4,5$ the upper bounds are four explicit rational inequalities.
For $t\ge6$, use $s_t\le3/2\le b^5\le b^{t-1}$.
The seed $[\beta+1](\rho)\le5/2$ and $b^\beta\ge1$ establish
(1) at the empty capacity word, including $\beta=0$.
If (1) holds before appending $t$, then $F'=1+s_tF$ and $M'=M+t-1$.
For $t\ge3$, the negative offset pays the whole new unit:
$(125/22)(s_t-1)\ge1$. For $t=2$, use $b^M\ge1$ and the exact
positive margin
\[
           10(b-24/25)-1-(125/22)/25=97/880>0.
\]
This proves the ENTIRE arbitrary-length fold, not a sampled range.

\subsection{A count-free first-missing penalty and a uniform power block}
Let $K=\kappa_Y\ge20$ and $n=K-2\ge18$. The LAST $n$ non-one
capacities of the ordered actual arm $Y$ are respectively at least
$2,3,\ldots,n+1$. Their skeleton degree is at least
$D_n=n(n+1)/2$. Write $M_0$ for the earlier prefix's skeleton degree.
Since $s_t\le c$, the whole tail fold satisfies
\[
 F_Y(\rho)+2\le c^n(F_{Y_0}(\rho)+2)
                  \le10c^n b^{M_0}.                  \tag{3}
\]
The last bound RETAINS the negative offset $125/22>2$ in (1).
The exact integer power checks are
\[
 q^{224}\le\tfrac12,\qquad
           2240c^{18}\le b^{171},\qquad c\le b^5.    \tag{4}
\]
They imply, for EVERY integer $m\ge0$ and $n\ge18$,
\[
                 m q^m\le224,\qquad2240c^n\le b^{D_n}. \tag{5}
\]
Indeed write $m=224a+r$, $0\le r<224$. Then
$m q^m\le224(a+1)2^{-a}\le224$ by $a+1\le2^a$.
For the second assertion start at $n=18$, use
$D_{n+1}=D_n+n+1$, and multiply by $b^{n+1}\ge b^5\ge c$.
Thus (5) is a whole unbounded block/triangle induction.

\subsection{The ACTUAL forward slope is dyadically thin outside its strict CI head}
Let $M=\beta+\sum(t-1)$ be the main-box degree of $Y$, and
$B=\max(\beta,\max(t-1))$. Put
\[
                            s_Y=\min(\lfloor M/2\rfloor,M-B).
\]
Then for EVERY $m>s_Y$,
\[
                          m\,\Delta F_Y(m)\le2^m.    \tag{6}
\]
First suppose $M<2m$. Positivity of coefficients, (3), and
$M\ge M_0+D_n$ give
\[
 b^{D_n}\frac{F_Y(m)}{2^m}
        \le10c^n q^m.
\]
Multiplication by $m$ and (5) prove $mF_Y(m)\le2^m$,
which is stronger than (6).

It remains to treat $m>M-B=:O$. The first-missing staircase gives
$O\ge1+(K-3)(K-2)/2\ge154$.
If the largest interval is the last capacity interval $[T]$, delete it
to obtain the actual prefix $Y_0$ with skeleton degree $O$. The exact
forward recurrence gives
\[
 \Delta F_Y(m)=F_{Y_0}(m-1)-F_{Y_0}(m-T-1)\le F_{Y_0}(m-1).
\]
Write $j=m-1\ge O$. From (1) and positivity,
\[
                   \frac{\Delta F_Y(m)}{2^m}\le5v^j.
\]
If instead the initial interval $[\beta+1]$ is largest, write
$F_Y=A+P[\beta+1]$ by positive suffix boxes, where
$A$ is their proper-suffix sum and $P=\prod S_t$.
Then coefficientwise
\[
 (1-z)F_Y=(1-z)A+P(1-z^{\beta+1})\preceq A+P=:F_{0,V}.
\]
This is the allowed zero-initial-run formal fold of skeleton degree $O$,
so (1) gives $\Delta F_Y(m)/2^m\le10v^m$.
The exact base $10\cdot101v^{100}<1$, followed by the decreasing-ratio
induction $v(j+2)\le j+1$ for $j\ge100$, proves
$10(j+1)v^j\le1$ for EVERY $j\ge100$. Since $O\ge154$,
both largest-factor cases prove (6). These cases cover the actual minimum
in $s_Y$, including ties, plateaus and zero/negative slopes.

\subsection{An indexed curvature gap for ANY actual mixed owner}
Let $Z=1^\alpha U$ have $h\ge2$ non-one parts, skeleton degree $D$,
weight $a=\alpha+D+h$, and ordinary reversal $Q=q_Z$.
At a mixed/source-prefix recovery rank let $r=a-j\ge0$,
$Q_r\ge Q_{r-1}$ and
\[
                           C=Q_{r+1}-2Q_r+Q_{r-1}>0.
\]
The frozen primitive-center theorem gives $D\ge2r+2$.
The frozen ordinary Hilbert realization gives
$(r+1)Q_{r+1}\le(h+r+1)Q_r$, including every padding and capacity.
We claim the stronger indexed gap
\[
                    (j+1)(2Q_r-Q_{r-1})\ge4Q_{r+1}.  \tag{7}
\]
For $r\ge2$ it suffices to show
$(j+1)(r+1)\ge4(h+r+1)$, because $2Q_r-Q_{r-1}\ge Q_r$.
We have $j+r\ge2h+1$ and $j\ge h+r+3$.
For $r\ge3$ use
$(h+r+4)(r+1)-4(h+r+1)=h(r-3)+r(r+1)\ge0$.
For $r=2$, split at $h=6$ and use the respective two bounds on $j$.
For $r=1$, use $Q_0=1,Q_1=h+1$,
$2Q_2\le(h+1)(h+2)$ and $j\ge2h$; the remaining positive polynomial
is $2h^2-2h-3\ge1$ at $h\ge2$.
For $r=0$, $Q_{-1}=0,Q_0=1,Q_1=h+1$ and $j\ge2h+1$ give (7)
directly. No first-missing bound or C assumption is used for the owner.

\subsection{A whole dyadic source sum pays the indexed gap}
Write $A=F_Z$, $B=F_Y$, $N=j+1$, and
$E_B=1-z+(2z-1)B$. Its entire prefix through $N$ is nondecreasing
in the mixed region; put $w_s=\Delta E_B(s)\ge0$.
Its first possible birth is at least three. The actual doubling law
$A_{i+1}\le2A_i$ gives, for every used $s$,
\[
                       A_{N-s}\ge2^{2-s}A_{N-2}.
\]
The COMPLETE telescoping identity, including the rank-two correction, is
\[
 4\sum_{s=2}^N2^{-s}w_s
                   =1-4\,2^{-N}\Delta B_N.           \tag{8}
\]
It follows from $B_0=B_1=1$ and $E_B(s)=2B_{s-1}-B_s$ for $s\ge2$.
Consequently
\[
 \Delta(E_BA)_N\ge A_{N-2}(1-4\,2^{-N}\Delta B_N).    \tag{9}
\]
The exact mixed tail increment is $T=C-\Delta(E_BA)_N$.
By (7), (9), and $N\Delta B_N\le2^N$ one obtains $T\le0$.
Thus positive recovery NECESSARILY has
\[
                           N\Delta F_Y(N)>2^N.         \tag{10}
\]
For $K\ge20$, (6) contradicts (10) whenever $N>s_Y$.
Therefore EVERY positive mixed recovery at nonnegative reverse rank has
\[
                              j+1\le s_Y.             \tag{11}
\]
This retains the actual opposite main degree and largest interval.
It does NOT replace them by a first-missing-only linear surrogate.

\subsection{All properly conditioned mixed CLOSED S heads at first missing twenty or above}
At both strict positive windows, the frozen reflection guard gives
$k+1\le j-2$ and $k+1\le a-1$. By (11), the ACTUAL signed opposite
differences satisfy $\Delta q_Y(i)\ge1$ throughout every used head:
the main CI has positive primitive support through $s_Y$, while each
proper shifted-CI summand has nonnegative differences there.
All used coefficients of the raw owner source $R_Z=(2-z)q_Z$ are
nonnegative by the second guard. If $0\le r<k$, its index $r+1$ term
pays $H\ge R_Z(r+1)\Delta q_Y(k-r)\ge C\ge T$.
The frozen relative raw-kernel theorem with $+1$ surplus pays $r\ge k$,
including CLOSED source equality $0\le E_n=E_k$.

For $r\le-2$ positive recovery is impossible. At $r=-1$ the first
source birth pays the unit curvature unless $a\le K-3$, whence
$j\le K-2$ and $k+1\le K-4\le s_Y$; the raw constant term pays.
One-part owners have only this same negative-rank birth exception.
For nonpositive $T$, the usual OLD-A2 head gives $H\ge0$.
This proves the ENTIRE mixed proper CLOSED S comparison for EVERY
$K\ge20$, without capacity, count, weight, padding or amplification bounds.

\section{The fourteen-terminal-part bound}

\noindent{\small Source\enspace\path{fixed_terminal_fourteen_slope_thinness_proof.tex}.}\par\medskip


Keep $\rho=3/5$, $b=35/32$, $q=6125/6144$, $v=175/192$,
$K_0=125/22$, and $s_t=S_t(\rho)$. Put $r_t=s_t/b^{t-1}$.
The frozen norm gives $F(\rho)+K_0\le10b^M$ for EVERY formal word.
For every append $t\ge3$, its stronger exact step is
\[
       F'(\rho)+K_0\le s_t(F(\rho)+K_0),              \tag{1}
\]
because $K_0(s_t-1)\ge1$. Thus an ENTIRE arbitrary tail of capacities
at least three has multiplicative transport of this offset norm.

\subsection{Three new sharp but elementary unbounded scalar bounds}
The ratios $r_t$ decrease for all $t\ge4$, and $r_3\le r_4$.
Indeed $s_{t+1}=s_t+\rho^{t+1}$ and, for $t\ge4$,
\[
 \rho^{t+1}\le\rho^5\le(b-1)s_4\le(b-1)s_t.
\]
Consequently $r_t\le r_{\max(4,L)}$ whenever $t\ge L\ge3$.
For ten arbitrary capacities $u_i\ge i+3$, $i=0,\ldots,9$,
\[
 \prod_{i=0}^9r_{u_i}\le P_*:=r_4^2\prod_{t=5}^{12}r_t.        \tag{2}
\]

A sharper WHOLE block gives
\[
                      m q^m\le128\qquad(m\ge0).       \tag{3}
\]
Use $q^{16}\le20/21$, write $m=16a+r$, $0\le r<16$, and
$(a+1)(20/21)^a\le8$. The latter sequence increases through index 19
and decreases thereafter; alternatively check indices 0..19 exactly
and use the decreasing-ratio induction after 19. This is a fixed
scalar base plus unbounded induction, with NO finite word premise.

For EVERY integer $i\ge0$,
\[
                       (i+1)\rho^i\le6/5.             \tag{4}
\]
Index zero is one, index one is $6/5$, and the consecutive ratio is
at most one after index one. Since $[u][w]$ has coefficient at rank
$i$ at most $i+1$, (4) gives, for every nonnegative polynomial $G$,
\[
              \rho^k[z^k]([u][w]G)\le(6/5)G(\rho).   \tag{5}
\]
This is a WHOLE convolution bound: bound each weighted kernel
coefficient and sum all used nonnegative $G$ terms. It does not
differentiate a coefficient inequality.

The single exact terminal-template constant is
\[
                \frac{41472}{125}P_*\le b^{39}.        \tag{6}
\]
The ratio of the two sides is approximately $0.9764430749$;
only the exact rational inequality, not this decimal, is a proof input.

\subsection{Bulk slope bound from a fixed last fourteen parts}
If an ACTUAL increasing arm has first missing weight at least sixteen,
it has at least fourteen non-one parts, and its LAST fourteen capacities
are at least $2,3,\ldots,15$. This fixed template suffices; no upper
bound on the first missing weight or any capacity is imposed.

Write the last three capacities as $w,u,T$, in nondecreasing order,
and let $G$ be the earlier forward arm, of skeleton degree $M_G$.
The first of the terminal fourteen capacities is absorbed into the
earlier prefix. The following TEN capacities are respectively at least
$3,4,\ldots,12$, so the frozen global norm and (1)--(2) give
\[
                       G(\rho)+K_0\le10b^{M_G}P_*.
                                                               \tag{7}
\]
The three deleted capacities are at least $13,14,15$, so their total
skeleton degree is at least $12+13+14=39$.

Consider any $m$ with $M\le2m$, where $M$ is the whole main-box degree.
The prefix before the LAST TWO parts has skeleton degree at least
$1+1+\cdots+12=79$. Since $T\ge u$, this implies $M\ge2u+4$.
Hence $m-1>u$. The exact forward identities therefore give
\[
 \begin{split}
 \Delta F_Y(m)&\le F_{Gwu}(m-1),\\
 F_{Gwu}&=1+S_u+S_uS_wG,\\
 F_{Gwu}(m-1)&=[z^{m-3}]([u][w]G).
 \end{split}                                                    \tag{8}
\]
Both source terms in the middle line vanish at the used coefficient;
this support guard is essential. Apply (5), (7), and the removed degree:
\[
 b^{39}\frac{\Delta F_Y(m)}{2^m}
     \le10\rho^3(6/5)q^m P_*.
\]
Multiplication by $m$ and (3)--(6), using
$10\rho^3(6/5)\cdot128=41472/125$, prove
\[
                          m\Delta F_Y(m)\le2^m.       \tag{9}
\]
No curvature, C, A1, A2 or source-prefix premise is used for this slope bound.

\subsection{Other-degree branch and all mixed proper heads}
Let $s_Y=\min(\lfloor M/2\rfloor,M-B)$, with $B$ the largest interval
degree. The nonbulk alternative $m>M-B$ is exactly the same as in the
frozen twenty-threshold proof. For first missing at least sixteen,
$M-B\ge1+13\cdot14/2=92$. Replace the old large-read base 100 by 80:
\[
                10(j+1)(175/192)^j\le1\qquad(j\ge80),
\]
whose base is an exact rational check and whose consecutive ratio
decreases thereafter. This covers BOTH a largest last capacity and a
largest initial interval, including the zero-initial-run formal comparison.
Combining with (9) gives EVERY $m>s_Y$ the signed slope bound (9).

The frozen COMPLETE indexed-owner gap and dyadic-prefix/source-sum
theorems then imply every positive mixed recovery has $j+1\le s_Y$.
All properly conditioned CLOSED S heads are paid by the genuine strict
opposite derivative prefix, the one-term absolute budget, the frozen
relative $+1$ surplus and the negative-birth unit argument. Nonpositive
recoveries retain the usual OLD nonnegative head budget.

\section{First missing weights thirteen through fifteen}

\noindent{\small Source\enspace\path{complete_mixed_source_unrestricted_C_original_criterion_proof.tex}.}\par\medskip

Keep $\rho=3/5$, $b=35/32$, $q=6125/6144<1$, $v=175/192$,
$K_0=125/22$, and $r_t=S_t(\rho)/b^{t-1}$. Frozen whole-fold norm:
$F(\rho)+K_0\le10b^M$. Frozen whole offset contraction for all
capacities at least three and the ratio envelope $r_t\le r_{\max(4,L)}$
for $t\ge L\ge3$ are available.

\subsection{A new constant-density theorem for EVERY first missing weight at least thirteen}
Theorem: for ANY actual $Y$ with $\kappa_Y\ge13$,
\[
 s_Y=\min(\lfloor M/2\rfloor,M-B),\quad m>s_Y
       \quad\Longrightarrow\quad20\,\Delta F_Y(m)\le2^m.       \tag{1}
\]
This is a CONSTANT density bound, weaker than $m\Delta F\le2^m$
when $m>20$, and is deliberately paired below with a stronger owner gap.

Such a $Y$ has LAST ELEVEN capacities at least $2,3,\ldots,12$.
Absorb the first into the unrestricted earlier prefix. The next TWO
capacities are at least three and four. Let $G$ be the prefix after
these first three terminal capacities, with skeleton degree $M_G$.
Offset contraction and the ratio envelope give
\[
                      G(\rho)\le10b^{M_G}r_4^2.       \tag{2}
\]
There remain EIGHT capacities at least $5,6,\ldots,12$, with total
skeleton degree at least $4+5+\cdots+11=60$.

Delete the LAST capacity $T$. Then
$\Delta F_Y(m)\le F_{Y_0}(m-1)$.
Expand the other SEVEN last capacities of $Y_0$:
\[
 F_{Y_0}=\sum_{\ell=0}^6\prod_{\text{last }\ell}S_t
             +\left(\prod_{\text{last }7}S_t\right)G.          \tag{3}
\]
The positive source sum at rank $m-1$ has at most the number of
compositions of $m-1$ with at most six parts. Their cut sets have size
at most five. List each cut set in increasing order and pad with zeroes
to a five-tuple; this is an injection into $m^5$ tuples. Thus its
coefficient $U_m$ satisfies $U_m\le m^5$ for $m\ge2$.
The empty-product source has zero coefficient at the ranks used here.

For the SEVEN interval kernel, stars and bars give
\[
 [z^i]\prod_{\text{last }7}[t]\le\binom{i+6}{6},\qquad
                \binom{i+6}{6}\rho^i\le3003\rho^8.     \tag{4}
\]
The last sequence has its two maxima at indices eight and nine.
Its consecutive ratio is $\rho(i+7)/(i+1)$; this gives a whole
unbounded proof, not a chosen coefficient range. The same WHOLE
positive weighted-convolution argument as the frozen two-interval
bound now yields
\[
 \rho^m\bigl(F_{Y_0}(m-1)-U_m\bigr)
                         \le3003\rho^{16}G(\rho).    \tag{5}
\]
It is ESSENTIAL to retain $U_m$; those seven source products do NOT
all vanish. Their entire contribution is paid separately.

In the BULK $M\le2m$, the eleven-part staircase gives $M\ge67$,
so $m\ge34$, and $M_G+60\le M\le2m$. Exact rational bases are
\[
 25\cdot30030\rho^{16}r_4^2\le b^{60},\qquad
              100\cdot34^5\le2^{34}.                 \tag{6}
\]
The whole induction $100m^5\le2^m$ holds for every $m\ge34$:
$(m+1)/m\le35/34$ and $(35/34)^5<2$.
Equations (2), (5), (6) and $q^m\le1$ give the long-kernel
density at most $1/25$. The ENTIRE source sum gives density at most
$m^5/2^m\le1/100$. Their sum is $1/20$, proving (1) in the bulk.
Numerically the long-kernel constant before allocating $1/25$ is
approximately $0.03899449166$; only the exact inequality (6) is used.

In the NONBULK $m>M-B$, the old largest-factor proof applies with
smaller degree thresholds. If the last capacity interval is largest,
$O=M-B\ge1+10\cdot11/2=56$, $j=m-1\ge O$, and
$\Delta F_Y(m)/2^m\le5v^j\le1/20$, from $100v^{56}\le1$.
If the initial interval is largest, its other degree is at least
$1+\cdots+11=66$, $m\ge67$, and
$\Delta F_Y(m)/2^m\le10v^m\le1/20$, from $200v^{67}\le1$.
Both are unbounded decreasing-power statements; zero and negative
slopes cause no issue. This proves (1) for all capacities, counts, weights and initial runs.

\subsection{The remaining first-missing values force a FIVE-to-one raw owner gap}
Take a positive mixed/source-prefix recovery with rising-opposite
first missing $K\in\{13,14,15\}$. At nonnegative owner reverse rank
$r$, frozen first-source and primitive facts give
\[
 D=2r+2+e,\quad h\le D,\quad \alpha+e+3\le K,\quad
       j=\alpha+h+r+2+e.                              \tag{7}
\]
If $j-2\le s_Y$, the genuine strict opposite head already pays every
proper absolute head. Otherwise
\[
 j-2>s_Y\ge S_K=\left\lfloor\frac{1+(K-2)(K-1)/2}{2}\right\rfloor,
 \qquad(S_{13},S_{14},S_{15})=(33,39,46).
\]
The finite THREE-value arithmetic of (7), without any finite owner
enumeration, implies
\[
                               e\le2r+2.              \tag{8}
\]
This is an exact unbounded linear-integer implication, checkable by
splitting ONLY $K=13,14,15$. In particular
$h+r+1\le3r+3+e\le5(r+1)$.
The frozen ordinary Hilbert growth and mixed owner phase now give,
with $Q=q_Z(r+1)$ and $R=2q_Z(r)-q_Z(r-1)$,
\[
 (r+1)Q\le(h+r+1)q_Z(r)\le5(r+1)q_Z(r),\quad
                   R\ge q_Z(r),\quad\boxed{Q\le5R}.  \tag{9}
\]

Put $N=j+1$. The frozen WHOLE dyadic source-prefix sum pays
\[
             \text{source}\ge Q\left(1-4\,2^{-N}\Delta F_Y(N)\right).
\]
Here $j-2>s_Y$ gives $N>s_Y$, so (1) implies
$4\,2^{-N}\Delta F_Y(N)\le1/5$.
By (9), the whole source pays $Q-R=C$, hence $T\le0$, contradicting
positive recovery. Thus every positive mixed recovery at $K=13,14,15$
is inside the genuine strict-head range. The negative-rank birth is
paid by the first source unless $j\le K-2$, also inside that head.
One-part owners have only this same birth exception.
Frozen absolute one-term and relative $+1$ head payments therefore
prove ENTIRE correctly conditioned MIXED proper CLOSED S at these
last three values. No arbitrary opposite capacity/count/padding cutoff
or finite word certificate is a premise.

\section{Proper deletion and the peak join}

\noindent{\small Source\enspace\path{inductive_signed_window.tex}.}\par\medskip

Write $[m]=1+\cdots+z^{m-1}$, $q_T=z^{\operatorname{wt}T}F_T(1/z)$,
and
\[
 \delta_M(T)=([M]-z^M)q_T+z^{\operatorname{wt}T+M}.
\]
This is the reversal of the nonnegative single-arm relative kernel
$1+(S_M-1)F_T$. In particular $\delta_M(\varnothing)=[M]$.
All sequences below have zero extension at negative indices.

\begin{lemma}[A proper relative polynomial supplies the signed prefix]
Let $X,Y$ be nonrectangular increasing arms below a strict peak $M$.
Orient their weights as $a=\operatorname{wt}X\le b=\operatorname{wt}Y$,
write $Y=Zt$ and $c=b-t$, and put
\[
 q=q_Xq_Y,\quad \sigma=\min\operatorname{argmax}q,
 \quad d=(\Delta q)^+.
\]
If $q_X\delta_M(Z)$ is unimodal, then $([M]-z^M)d$
has no negative-to-positive coefficient change through rank $\sigma$.
\end{lemma}
\begin{proof}
Ordinary two-arm unimodality and its strengthened midpoint bound
give $\sigma<b$. The strict peak gives $c+M=b+(M-t)>b$.
The exact reverse append identity is
\[
 q=[t]q_Xq_Z+z^bq_X.
\]
Set $V=[t]q_X\delta_M(Z)$. Then
\[
 V=([M]-z^M)q-z^b([M]-z^M)q_X
                      +z^{c+M}[t]q_X.
\]
The two corrections start strictly after $\sigma$. Therefore
$\Delta V_j=([M]-z^M)\Delta q_j$ for every $0\le j\le\sigma$.
On this prefix $q$ is nondecreasing, so $d_j=\Delta q_j$;
causality of multiplication by a polynomial with nonnegative powers
gives
\[
                 [z^j]([M]-z^M)d=\Delta V_j
                       \qquad(0\le j\le\sigma).
\]
Uniform convolution preserves unimodality, so the premise makes
$V=[t](q_X\delta_M(Z))$ unimodal. Its differences cannot have
a negative coefficient followed by a positive one. This proves
the claimed prefix property. No assertion about the later
coefficients of $([M]-z^M)d$ is needed.
\end{proof}

\begin{corollary}[Only the A2 premise is needed in a least-counterexample proof]
At a least relative counterexample at fixed peak $M$, the signed
prefix premise in {\normalfont\path{general_base_join_route.tex}} holds
automatically. Consequently the general unamplified one-descent
conjecture would follow from the single assertion
\[
       2^j\max(\Delta(q_Xq_Y)_j,0)\quad\hbox{is unimodal}
       \tag{A2}
\]
for every pair of nonrectangular increasing arms.
\end{corollary}
\begin{proof}
Choose the least failing peak, then a relative counterexample with
the fewest total arm parts. Empty or rectangular arms are already
known. The established relative deletion reduction gives
$|a-b|\le M-4$ and equivalence of core and both relative
unimodality in this band. The relative polynomial with arms
$X,Z$ has fewer parts, in the required orientation as well,
so it is unimodal. Its reversal is $q_X\delta_M(Z)$; the lemma
supplies the signed prefix premise.

Take $t_0$ to be the largest arm part. The base polynomial
$J_{t_0}=A+B+(S_{t_0}-1)AB$ is unimodal: remove all arm
occurrences of $t_0$ using $J=1+S_{t_0}J^{\rm smaller}$,
and apply the least-peak hypothesis to the remaining strict
peak $t_0<M$ (or the elementary empty-arm case).
The conditional general base-peak join, now using only (A2)
as a new structural input, makes the core unimodal. Balanced
core-relative equivalence contradicts the chosen counterexample.
\end{proof}

\noindent{\small Source\enspace\path{short_prefix_reflection_join.tex}.}\par\medskip

Let $a\le b$ be the weights of two nonrectangular increasing arms,
$D=a+b$, $N=D+M$, $R=\lfloor N/2\rfloor$, and
\[
 q=q_Xq_Y,\quad \sigma=\min\operatorname{argmax}q,
 \quad d_j=\max(q_j-q_{j-1},0).
\]
Established ordinary-product results give unimodality of $q$,
unimodality of $d$, $d_0=1$, $\sigma<b$, and reverse reflection
\[
                         q_{D-n}\ge q_n\qquad(n\ge D/2).
\]
The last inequality follows by reversing the ordinary height-layer
endpoint sums, each at least $D$.

\begin{lemma}[A finite halving tail with one spare endpoint]
Suppose $d\ge0$ is unimodal and $d_0>0$.
If $k\le a-2$ and $2d_{i+1}\le d_i$ for $k\le i\le a-2$,
then
\[
                   \sum_{i=k+1}^{h}d_i\le d_k
                              \qquad(k\le h\le a).
\]
\end{lemma}
\begin{proof}
For $h\le a-1$, sum $d_{i+1}\le d_i-d_{i+1}$ to get
the stronger bound $\sum_{i=k+1}^h d_i\le d_k-d_h$.
The first halving edge forces the global unweighted mode to
be at most $k$: otherwise unimodality and $d_0>0$ would give
$0<d_k\le d_{k+1}\le d_k/2$, a contradiction.
Thus $d_a\le d_{a-1}$. This pays the one additional endpoint
when $h=a$.
\end{proof}

\begin{proposition}[A shorter-prefix base-peak join]
Assume $b>M$ and $b-a\le M-4$. Suppose
\begin{enumerate}
\item the finite sequence $(2^jd_j)_{0\le j<a}$ is unimodal;
\item the proper relative polynomial $q_X\delta_M(Y^-)$ is unimodal;
\item for some $2\le t<M$, the base
$J_t=A+B+(S_t-1)AB$ is unimodal.
\end{enumerate}
Then the strict core $H_M$ is unimodal.
\end{proposition}
\begin{proof}
Put $w=([M]-z^M)q$. Before rank $R+1$ the reverse core
equals $w$, because the arm corrections start at $a+M,b+M>R$.
The exact proper-relative identity proved in
{\normalfont\path{inductive_signed_window.tex}} shows that $w$ agrees,
through rank $b-1$, with the unimodal polynomial
$V=[\max Y]q_X\delta_M(Y^-)$.
Thus $\Delta w$ has no negative-to-positive change there.

Choose a mode $\ell$ of the truncated weighted sequence in
the first premise. It may be chosen with
$\ell\le\min(\sigma,a-1)$. For
\[
              b\le n=M+k\le\min(R,M+\sigma)
\]
one has $0\le k\le a-2$ and the exact difference is
\[
                   J_k:=\Delta w_{M+k}=q_{M+k}-q_k-d_k.
\]
If $k<\ell$, weighted increase gives
$2d_{k+1}\ge d_k$. Hence $q_k+d_k$ is nondecreasing,
while $q_{M+k}$ is nonincreasing since $M+k\ge b>\sigma$.
The $J_k$ are therefore nonincreasing in this range. The same
comparison includes the joining edge from $n=b-1$ whenever
the new range begins before $M+\ell$; here $b-1\ge\sigma$.

If $k\ge\ell$, set $h=D-M-k=D-n$. The bounds on $n$
give $k\le h\le a$, and reflection gives $q_n\le q_h$.
Weighted decrease gives every halving edge through $a-1$.
The preceding lemma yields
\[
 q_n-q_k\le q_h-q_k
       \le\sum_{i=k+1}^{h}d_i\le d_k.
\]
Thus $J_k\le0$. These comparisons extend the no-revival
prefix through $\min(R,M+\sigma)$.

For the remaining ranks $M+\sigma<n\le R$, use the established
base-peak join exactly as in {\normalfont\path{general_base_join_route.tex}}:
the last forward product mode $p=D-\sigma$ exceeds both arm weights,
$\Delta J_t(p+1)>0$, and
$H_M=J_t+z^{t+1}[M-t]AB$ supplies forward increase at each
corresponding rank. Thus all these remaining reverse differences
are nonpositive. Forward first-half growth joins the final reverse
tail, proving core unimodality.
\end{proof}

The base-peak step remains necessary in this argument: past
$M+\sigma$, $\Delta q_k$ is negative and cannot be replaced by
its positive part in the displayed formula for $J_k$.

\begin{corollary}[A still weaker universal target]
The general conjecture follows if, for every nonrectangular
arm pair, $(2^jd_j)_{0\le j<\min(a,b)}$ is unimodal.
It is enough instead to prove coefficient unimodality of
$(2-z)q_Xq_Y$ restricted to this prefix. The full A2 theorem
or full mixed-deficit theorem is not required by this route.
\end{corollary}
\begin{proof}
Use a least-peak, least-relative counterexample. Known small
peaks and $N\le3M+8$ are excluded, so $b>M$; balance gives
the other size hypothesis. Proper-relative deletion gives the
second premise, and least-peak induction supplies $J_t$ after
stripping repeated arm maxima.
For the alternative premise, through the first reverse mode,
\[
 \Delta((2-z)q)_j=2d_j-d_{j-1}.
\]
These are, up to a positive factor, the successive differences
of $2^jd_j$. Truncate at the earlier of the prefix endpoint
and $\sigma$; subsequent positive differences of $q$ are zero.
Hence prefix unimodality of $(2-z)q$ supplies the required
truncated weighted-difference property.
\end{proof}

\noindent{\small Source\enspace\path{general_base_join_route.tex}.}\par\medskip

\begin{proposition}[Conditional general propagation]
Assume these two properties for the given increasing arms and
strict peak $M$. If $2\le t<M$ and
$J_t=A+B+(S_t-1)AB$ is unimodal, then the balanced core
$H_M=A+B+(S_M-1)AB$ is unimodal whenever $|\deg A-\deg B|<M$.
\end{proposition}
\begin{proof}
Put $w=([M]-z^M)q$ and choose a mode $\ell$ of $2^jd_j$.
Through rank $s$, its difference sequence agrees with
$([M]-z^M)d$, so the second premise gives single crossing.
For $n=M+k$ in the crossing range $k+M\ge s$, $k<\ell$,
the first premise gives $d_{k+1}\ge d_k/2$. Therefore
$q_k+d_k$ is nondecreasing and $q_{k+M}$ is nonincreasing,
making $\Delta w_{M+k}=q_{k+M}-q_k-d_k$ nonincreasing.
For $\ell\le k<s$, the same premise instead gives
\[
 \Delta w_{M+k}\le d_k\left(\sum_{i=1}^M2^{-i}-1\right)
                    =-2^{-M}d_k\le0.
\]
At $k=s$ this difference is again nonpositive. Thus the reverse
core has the needed single crossing through rank $M+s$.

The rest is the same base-peak join as in the binary theorem.
If a remaining reverse rank $M+s<n\le\lfloor N/2\rfloor$
exists, then $p=D-s>\max(\deg A,\deg B)$. Hence
\[
 \Delta J_t(p+1)=2P_p-P_{p+1}-P_{p-t}>0.
\]
Its unimodality gives growth through $p+1$. The identity
\[
                      H_M=J_t+z^{t+1}[M-t]P
\]
then gives $\Delta H_M(j+1)=\Delta J_t(j+1)+P_{j-t}-P_{j-M}
\ge0$ at every corresponding $j=N-n<p$.
Forward first-half growth completes the proof.
\end{proof}

\end{document}
